% !TeX program = xelatex
% !TeX encoding = UTF-8
% !TeX spellcheck = pt_BR
% ----------------------------------------------------------------------
% Guia Prático para a Capacitação de Git e GitHub
% Compilar com XeLaTeX (a fonte Lexend é carregada localmente):
%     latexmk -xelatex guia.tex
% ----------------------------------------------------------------------

\documentclass[12pt, a4paper, openright, oneside, brazil]{abntex2}

% ----------------------------------------------------------------------
% Preâmbulo — Capacitação de Git e GitHub
% ----------------------------------------------------------------------
% ATENÇÃO: compilar com XeLaTeX. A fonte Lexend é carregada de arquivos
% locais em assets/fonts/, o que exige fontspec (XeLaTeX ou LuaLaTeX).
% ----------------------------------------------------------------------
\usepackage[brazil]{babel}

% --- Fonte da marca (Lexend, local) -----------------------------------
\usepackage{fontspec}
\newcommand{\fontpath}{./assets/fonts/}
\setmainfont[
    Path          = \fontpath,
    UprightFont   = *-Regular,
    BoldFont      = *-Bold,
    AutoFakeSlant = 0.2,          % Lexend não possui itálico real
]{Lexend}

% --- Estrutura e tipografia -------------------------------------------
\usepackage{microtype}
\usepackage{csquotes}      % \enquote{...}: aspas no padrão brasileiro
\usepackage{indentfirst}
\usepackage{graphicx}
\usepackage{xcolor}
\usepackage{float}
\usepackage{fancyhdr}

% --- Tabelas -----------------------------------------------------------
\usepackage{longtable}
\usepackage{tabularx}
\usepackage{booktabs}
\usepackage{array}
\usepackage{multirow}

% --- Código ------------------------------------------------------------
\usepackage{listings}

% --- Símbolos usados nos avisos ---------------------------------------
\usepackage{amssymb}
\usepackage{pifont}

% --- Caixas de destaque ------------------------------------------------
\usepackage[most]{tcolorbox}

% --- Links e citações --------------------------------------------------
\usepackage{url}
\usepackage{hyperref}
\usepackage{abntex2cite}

% --- Justificação -----------------------------------------------------
% Nomes de comandos em \texttt{} são longos e não hifenizam, o que empurra
% linhas para fora da margem. O emergencystretch autoriza o TeX a esticar
% mais os espaços nesses parágrafos em vez de estourar a caixa.
\setlength{\emergencystretch}{3em}

% ----------------------------------------------------------------------
% Configurações visuais — Capacitação de Git e GitHub
% ----------------------------------------------------------------------

% --- Paleta da marca ---------------------------------------------------
\definecolor{branco}{HTML}{FFFFFF}
\definecolor{roxoescuro}{HTML}{1E1647}
\definecolor{roxoclaro}{HTML}{614AD3}
\definecolor{amarelo}{HTML}{FFFF00}
\definecolor{cinzafundo}{HTML}{F4F4F7}
\definecolor{cinzalinha}{HTML}{9A9AB0}
\definecolor{verdeok}{HTML}{1B7F4B}
\definecolor{vermelhoerro}{HTML}{B3261E}

% --- Atalhos de marcação -----------------------------------------------
\newcommand{\comando}[1]{\texttt{#1}}
\newcommand{\sqlkw}[1]{\texttt{\textbf{#1}}}
\newcommand{\arq}[1]{\texttt{#1}}
\newcommand{\menu}[1]{\textbf{#1}}
\newcommand{\tecla}[1]{\fbox{\footnotesize\ttfamily #1}}
\newcommand{\ok}{\textcolor{verdeok}{\ding{51}}}
\newcommand{\nok}{\textcolor{vermelhoerro}{\ding{55}}}

% ----------------------------------------------------------------------
% Ambientes de código (listings)
% ----------------------------------------------------------------------
% O 'literate' abaixo é necessário porque o listings não processa UTF-8
% sozinho: sem ele, cada acento dentro de um bloco de código vira lixo.
\lstset{
    inputencoding=utf8,
    extendedchars=true,
    literate=%
      {á}{{\'a}}1 {à}{{\`a}}1 {ã}{{\~a}}1 {â}{{\^a}}1
      {é}{{\'e}}1 {ê}{{\^e}}1 {í}{{\'i}}1
      {ó}{{\'o}}1 {õ}{{\~o}}1 {ô}{{\^o}}1
      {ú}{{\'u}}1 {ü}{{\"u}}1 {ç}{{\c c}}1
      {Á}{{\'A}}1 {À}{{\`A}}1 {Ã}{{\~A}}1 {Â}{{\^A}}1
      {É}{{\'E}}1 {Ê}{{\^E}}1 {Í}{{\'I}}1
      {Ó}{{\'O}}1 {Õ}{{\~O}}1 {Ô}{{\^O}}1
      {Ú}{{\'U}}1 {Ç}{{\c C}}1
      {º}{{\textordmasculine}}1 {ª}{{\textordfeminine}}1
}

\lstdefinestyle{base}{
    basicstyle       = \ttfamily\footnotesize,
    backgroundcolor  = \color{cinzafundo},
    frame            = single,
    rulecolor        = \color{cinzalinha},
    breaklines       = true,
    breakatwhitespace= false,
    postbreak        = \mbox{\textcolor{roxoclaro}{$\hookrightarrow$}\space},
    showstringspaces = false,
    upquote          = true,
    tabsize          = 2,
    captionpos       = b,
    xleftmargin      = 12pt,
    framexleftmargin = 12pt,
}

\lstdefinestyle{sql}{
    style        = base,
    language     = SQL,
    morekeywords = {AUTO_INCREMENT,ENGINE,InnoDB,UNSIGNED,IF,EXISTS,DATABASE,
                    REFERENCES,CONSTRAINT,CASCADE,RESTRICT,BOOLEAN,TINYINT,
                    DECIMAL,VARCHAR,TIMESTAMP,CHARACTER,COLLATE,LIMIT,SHOW,
                    DESCRIBE,TRUNCATE,USE,REPLACE},
    keywordstyle = \color{roxoescuro}\bfseries,
    commentstyle = \color{cinzalinha}\itshape,
    stringstyle  = \color{roxoclaro},
}

\lstdefinestyle{bash}{
    style        = base,
    language     = bash,
    keywordstyle = \color{roxoescuro}\bfseries,
    commentstyle = \color{cinzalinha}\itshape,
    stringstyle  = \color{roxoclaro},
}

\lstdefinestyle{yaml}{
    style         = base,
    basicstyle    = \ttfamily\footnotesize\color{roxoescuro},
    commentstyle  = \color{cinzalinha}\itshape,
    morecomment   = [l]{\#},
    moredelim     = [l][\color{black}]{:\ },
}

\lstdefinestyle{saida}{
    style           = base,
    basicstyle      = \ttfamily\footnotesize\color{roxoescuro},
    backgroundcolor = \color{branco},
}

\lstnewenvironment{sqlcode}[1][]{\lstset{style=sql,#1}}{}
\lstnewenvironment{saida}[1][]{\lstset{style=saida,#1}}{}
\lstnewenvironment{bashcode}[1][]{\lstset{style=bash,#1}}{}
\lstnewenvironment{yamlcode}[1][]{\lstset{style=yaml,#1}}{}

% ----------------------------------------------------------------------
% Caixas de destaque
% ----------------------------------------------------------------------
\newtcolorbox{aviso}[1][Atenção]{
    enhanced, breakable,
    colback   = vermelhoerro!5,
    colframe  = vermelhoerro,
    coltitle  = branco,
    fonttitle = \bfseries,
    title     = {#1},
    left=3mm, right=3mm, top=2mm, bottom=2mm,
}

\newtcolorbox{dica}[1][Dica]{
    enhanced, breakable,
    colback   = roxoclaro!5,
    colframe  = roxoclaro,
    coltitle  = branco,
    fonttitle = \bfseries,
    title     = {#1},
    left=3mm, right=3mm, top=2mm, bottom=2mm,
}

\newtcolorbox{conceito}[1][Conceito]{
    enhanced, breakable,
    colback   = roxoescuro!5,
    colframe  = roxoescuro,
    coltitle  = branco,
    fonttitle = \bfseries,
    title     = {#1},
    left=3mm, right=3mm, top=2mm, bottom=2mm,
}

% ----------------------------------------------------------------------
% Links e metadados do PDF
% ----------------------------------------------------------------------
\hypersetup{
    colorlinks  = true,
    linkcolor   = roxoescuro,
    citecolor   = roxoescuro,
    filecolor   = roxoclaro,
    urlcolor    = roxoclaro,
    pdftitle    = {Guia Prático para a Capacitação de Git e GitHub},
    pdfauthor   = {Ronivaldo Domingues de Andrade},
    pdfsubject  = {Git, GitHub, GitHub CLI, Git LFS, GitHub Pages e GitHub Actions},
    pdfkeywords = {Git, GitHub, GitHub CLI, Git LFS, GitHub Pages,
                   GitHub Actions, Pull Request, Assinatura de commits},
}


% --- Informações da capa e da folha de rosto --------------------------
\titulo{Git e GitHub na Prática: do Primeiro Commit à Automação}
\autor{Ronivaldo Domingues de Andrade}
\local{Rio de Janeiro - RJ}
\data{Edição 2025 · Revisão 2026}
\tipotrabalho{Guia Prático para Capacitação}
\preambulo{Guia prático da capacitação de Git e GitHub, com duração de duas
horas, elaborado para os membros da liga acadêmica. Cobre conta e benefícios
de estudante, instalação no Windows, no Linux e no macOS, Git, GitHub CLI,
Git LFS, commits, merges, Pull Requests, GitHub Pages, GitHub Actions e
assinatura de commits.}

\begin{document}
\selectlanguage{brazil}

% --- Elementos pré-textuais -------------------------------------------
\imprimircapa
\imprimirfolhaderosto

\newpage
\listoffigures

\newpage
\listoftables

\begin{siglas}
  \item[2FA] \emph{Two-Factor Authentication}, autenticação em dois fatores
  \item[CI/CD] \emph{Continuous Integration / Continuous Delivery}, integração e entrega contínuas
  \item[CLI] \emph{Command Line Interface}, interface de linha de comando
  \item[GCM] \emph{Git Credential Manager}, gerenciador de credenciais do Git
  \item[GPG] \emph{GNU Privacy Guard}, programa de criptografia e assinatura digital
  \item[LFS] \emph{Large File Storage}, extensão do Git para arquivos grandes
  \item[PAT] \emph{Personal Access Token}, token de acesso pessoal
  \item[PR] \emph{Pull Request}, pedido de integração de uma branch em outra
  \item[SSH] \emph{Secure Shell}, protocolo de acesso seguro por par de chaves
  \item[YAML] \emph{YAML Ain't Markup Language}, formato dos arquivos de workflow
\end{siglas}

\tableofcontents

% --- Capítulos ---------------------------------------------------------
% ======================================================================
\chapter{Introdução}
\label{cap:introducao}
% ======================================================================

O GitHub é hoje a principal plataforma de desenvolvimento colaborativo. Ele
guarda o código, registra quem mudou o quê, organiza a revisão entre colegas e
ainda publica sites e roda automações. Tudo isso em volta de uma ferramenta que
vem antes dele: o \textbf{Git}.

Este guia acompanha a capacitação de Git e GitHub da liga acadêmica. A
apresentação dura duas horas; aqui está tudo o que foi dito com mais calma, os
passos completos para Windows, Linux e macOS, e os exercícios para praticar
depois.

\section{Objetivos de aprendizagem}

Ao final, o participante é capaz de:

\begin{enumerate}
  \item criar e proteger a conta no GitHub, com autenticação em dois fatores, e
        pedir os benefícios de estudante;
  \item instalar e configurar o Git e o GitHub CLI no próprio sistema;
  \item registrar alterações com commits bem escritos e trabalhar com branches;
  \item integrar trabalho com merge e resolver um conflito;
  \item contribuir com um projeto de outra pessoa por fork e Pull Request;
  \item publicar um site com o GitHub Pages e automatizar uma tarefa com o
        GitHub Actions;
  \item versionar arquivos grandes com o Git LFS e assinar commits.
\end{enumerate}

\section{O roteiro das duas horas}

\begin{table}[H]
  \centering
  \caption{Roteiro de tempo da capacitação}
  \label{tab:roteiro}
  \begin{tabularx}{\textwidth}{@{}c X c c@{}}
    \toprule
    \textbf{Bloco} & \textbf{Conteúdo} & \textbf{Início} & \textbf{Duração} \\
    \midrule
    0 & Abertura, objetivos e combinados & 00:00 & 5 min \\
    1 & A conta no GitHub, 2FA e benefícios de estudante & 00:05 & 15 min \\
    2 & Instalação, configuração e autenticação & 00:20 & 20 min \\
    3 & Git no dia a dia: commits, branches, merge e conflito & 00:40 & 20 min \\
    — & \textit{Pausa} & 01:00 & 5 min \\
    4 & Colaboração: fork, Pull Request e padrões de commit & 01:05 & 20 min \\
    5 & GitHub Pages e GitHub Actions & 01:25 & 15 min \\
    6 & Git LFS e assinatura de commits & 01:40 & 15 min \\
    7 & Encerramento e próximos passos & 01:55 & 5 min \\
    \midrule
      & \textbf{Total} & & \textbf{120 min} \\
    \bottomrule
  \end{tabularx}
\end{table}

\section{Como ler este guia}

Os comandos aparecem em blocos como este. Linhas que começam com \texttt{\#}
são comentários e não precisam ser digitadas:

\begin{bashcode}
# confere a versão do Git instalada
git --version
\end{bashcode}

Ao longo do texto, três tipos de caixa chamam a atenção:

\begin{conceito}
Uma definição importante, para guardar.
\end{conceito}

\begin{dica}
Um atalho ou uma boa prática.
\end{dica}

\begin{aviso}
Um erro comum, ou um cuidado que evita dor de cabeça.
\end{aviso}

Quando um comando traz algo entre \texttt{<sinais>}, como
\texttt{<seu-usuario>}, troque tudo, inclusive os sinais, pelo valor real.

\begin{dica}[Qual terminal usar]
No \textbf{Windows}, use o \textbf{Git Bash}, que vem com o Git (capítulo
\ref{cap:ambiente}). Os comandos deste guia, como \texttt{ls}, \texttt{touch}
e \texttt{mkdir -p}, são do Bash e não funcionam no Prompt de Comando. No
\textbf{Linux} e no \textbf{macOS}, use o Terminal do sistema. A única exceção
é a instalação pelo \texttt{winget}, feita no PowerShell.
\end{dica}

\section{Pré-requisitos}

\begin{itemize}
  \item Um computador com Windows 10 (versão 1809 ou mais nova), Windows 11,
        Linux ou macOS, com permissão para instalar programas.
  \item Um endereço de e-mail. Para os benefícios de estudante, é melhor ter o
        e-mail acadêmico.
  \item Um celular com um aplicativo autenticador (Google Authenticator,
        Microsoft Authenticator ou similar), para a autenticação em dois
        fatores.
\end{itemize}

Não é preciso saber programar. Os exercícios usam arquivos prontos, que estão
na pasta \arq{materials/} do repositório.
          % Introdução e roteiro
% ======================================================================
\chapter{A conta no GitHub}
\label{cap:github}
% ======================================================================

\section{O que é o GitHub}

\begin{conceito}[Git e GitHub não são a mesma coisa]
O \textbf{Git} é um programa de controle de versão que roda no seu computador.
O \textbf{GitHub} é uma plataforma na internet que hospeda repositórios Git e
acrescenta o que uma equipe precisa para colaborar: revisão de código (Pull
Requests), controle de acesso, Issues, publicação de sites (GitHub Pages) e
automação (GitHub Actions).
\end{conceito}

O Git funciona sozinho, sem internet. O GitHub entra quando é preciso
compartilhar o trabalho, revisar o de outras pessoas ou publicar algo.

\section{Criando a conta}

\begin{enumerate}
  \item Acesse \url{https://github.com} e clique em \menu{Sign up}.
  \item Informe e-mail, senha e nome de usuário
        (Figura~\ref{fig:signup_page}). O nome de usuário aparece em todos os
        seus endereços, como \texttt{github.com/seu-usuario}: escolha um que
        você use profissionalmente.
  \begin{figure}[H]
    \centering
    \includegraphics[width=0.8\textwidth]{./assets/images/01_signup_page.png}
    \caption{Página de criação de conta no GitHub}
    \label{fig:signup_page}
  \end{figure}
  \item Confirme o e-mail pelo código enviado e conclua as telas seguintes. Ao
        final, você verá a página inicial (Figura~\ref{fig:signed_page}).
  \begin{figure}[H]
    \centering
    \includegraphics[width=0.8\textwidth]{./assets/images/02_signed_page.png}
    \caption{Primeira visão do GitHub depois de criar a conta}
    \label{fig:signed_page}
  \end{figure}
\end{enumerate}

\section{Autenticação em dois fatores (2FA)}

O GitHub exige a autenticação em dois fatores de quem contribui com código na
plataforma. Quando a exigência chega à sua conta, você recebe um aviso e tem
45 dias para ativar; depois disso, há mais 7 dias de carência antes de o
acesso ser bloqueado \cite{github_2fa}. Como todo participante desta
capacitação vai publicar código, vale ativar logo no primeiro dia.

\begin{enumerate}
  \item Clique na sua foto (canto superior direito) \(\rightarrow\)
        \menu{Settings} \(\rightarrow\) \menu{Password and authentication}.
  \item Em \menu{Two-factor authentication}, clique em
        \menu{Enable two-factor authentication}.
  \item Leia o QR code com um aplicativo autenticador no celular e digite o
        código de seis dígitos que ele mostrar.
  \item \textbf{Baixe e guarde os códigos de recuperação}. Eles são a única
        forma de entrar na conta se você perder o celular.
\end{enumerate}

\begin{dica}[Método principal e reserva]
O GitHub recomenda um \textbf{aplicativo autenticador} (TOTP) como método
principal e uma \emph{passkey} ou chave de segurança como reserva. Passkeys
ainda não são aceitas como método principal.
\end{dica}

\section{Benefícios para estudantes}

\subsection{O que o GitHub Education oferece}

Estudantes verificados pelo GitHub Education recebem, sem custo e enquanto
forem estudantes:

\begin{itemize}
  \item o plano \textbf{GitHub Pro} na conta pessoal;
  \item o \textbf{Copilot Student}, um plano do GitHub Copilot próprio para
        estudantes, ativado em um passo separado \cite{github_copilot_student};
  \item o \textbf{Student Developer Pack}, com ofertas de ferramentas de
        empresas parceiras.
\end{itemize}

\begin{table}[H]
  \centering
  \caption{Diferenças entre o GitHub Free e o GitHub Pro (conta pessoal)}
  \label{tab:free_pro}
  \begin{tabularx}{\textwidth}{@{}X c c@{}}
    \toprule
    \textbf{Recurso} & \textbf{GitHub Free} & \textbf{GitHub Pro} \\
    \midrule
    Repositórios públicos e privados & ilimitados & ilimitados \\
    Minutos de GitHub Actions por mês\textsuperscript{a} & 2.000 & 3.000 \\
    Armazenamento do GitHub Packages & 500 MB & 2 GB \\
    Codespaces: horas de núcleo por mês & 120 & 180 \\
    Codespaces: armazenamento por mês & 15 GB & 20 GB \\
    Git LFS: armazenamento e banda por mês & 10 GiB & 10 GiB \\
    GitHub Pages em repositório privado & \nok & \ok \\
    Revisores obrigatórios em repositório privado & \nok & \ok \\
    Suporte & comunidade & e-mail \\
    \bottomrule
  \end{tabularx}
  \par\smallskip
  {\footnotesize \textsuperscript{a} Os minutos só são contados em
  repositórios privados: em repositórios públicos, o GitHub Actions é gratuito.}
\end{table}

\begin{dica}[Copilot também existe no plano gratuito]
Qualquer conta pode usar o \textbf{Copilot Free}, com limites mensais (por
exemplo, 2.000 sugestões de código por mês). O Copilot Student libera mais
recursos para quem é estudante verificado \cite{github_copilot_plans}.
\end{dica}

\subsection{Quem pode pedir}

Segundo a documentação oficial \cite{github_education_apply}, é preciso:

\begin{itemize}
  \item estar matriculado em um curso que confere diploma (escola, faculdade,
        universidade);
  \item ter pelo menos 13 anos;
  \item ter uma conta pessoal no GitHub;
  \item enviar um comprovante de matrícula atual: foto da carteira de
        estudante com a data de validade, grade de horários, histórico ou
        declaração de matrícula.
\end{itemize}

\subsection{Passo a passo}

\begin{enumerate}
  \item \textbf{Cadastre o e-mail acadêmico.} Em \menu{Settings}
        \(\rightarrow\) \menu{Emails}, adicione o e-mail da instituição
        (Figura~\ref{fig:adding_academic_email}) e confirme pelo link que
        chegar na caixa de entrada.
  \begin{figure}[H]
    \centering
    \includegraphics[width=0.9\textwidth]{./assets/images/03_adding_academic_email.png}
    \caption{Adicionando o e-mail acadêmico}
    \label{fig:adding_academic_email}
  \end{figure}
  \item \textbf{Abra os benefícios de educação.} Acesse
        \url{https://github.com/settings/education/benefits} e, em
        \menu{GitHub Education}, clique em \menu{Start an application}
        (Figura~\ref{fig:github_student_developer_pack}). A página pública do
        pack (Figura~\ref{fig:github_student}) leva ao mesmo formulário.
  \begin{figure}[H]
    \centering
    \includegraphics[width=0.9\textwidth]{./assets/images/04_gsdp_page.png}
    \caption{Página do Student Developer Pack (\texttt{education.github.com/pack})}
    \label{fig:github_student}
  \end{figure}
  \begin{figure}[H]
    \centering
    \includegraphics[width=0.9\textwidth]{./assets/images/05_application_page.png}
    \caption{\menu{Education benefits} nas configurações, com o botão \menu{Start an application}}
    \label{fig:github_student_developer_pack}
  \end{figure}
  \begin{figure}[H]
    \centering
    \includegraphics[width=0.9\textwidth]{./assets/images/06_start_application.png}
    \caption{Formulário da candidatura}
    \label{fig:application_start}
  \end{figure}
  \item \textbf{Preencha e envie.} Escolha \menu{Student}, informe a
        instituição e o e-mail acadêmico (Figura~\ref{fig:application_start}),
        anexe o comprovante
        e clique em \menu{Submit application}. A análise pode levar alguns
        dias.
  \item \textbf{Ative o Copilot Student depois da aprovação.} A aprovação e a
        ativação do Copilot são passos separados
        \cite{github_copilot_student}. Volte a
        \url{https://github.com/settings/education/benefits} e siga as
        instruções de ativação. Se só aparecerem planos pagos, espere alguns
        dias: o benefício pode demorar para ser aplicado.
\end{enumerate}

\begin{dica}[Experiência do autor]
A foto da carteira de estudante, com nome, foto e data de validade visíveis,
costuma ser analisada mais rápido que outros documentos.
\end{dica}

\begin{aviso}[As telas mudam]
As capturas deste capítulo são de 2025. O GitHub reorganiza as telas com
frequência; se algo estiver diferente, procure pelos nomes dos botões citados
no texto.
\end{aviso}
          % A conta no GitHub e os benefícios de estudante
% ======================================================================
\chapter{Preparando o ambiente}
\label{cap:ambiente}
% ======================================================================

Duas ferramentas acompanham toda a capacitação:

\begin{itemize}
  \item o \textbf{Git}, que controla as versões no seu computador;
  \item o \textbf{GitHub CLI} (\texttt{gh}), que opera o GitHub pelo terminal:
        cria repositórios, abre Pull Requests e faz o login.
\end{itemize}

Siga só a seção do seu sistema.

% ----------------------------------------------------------------------
\section{Windows}
% ----------------------------------------------------------------------

No Windows, a instalação é feita pelo \textbf{winget}, o gerenciador de pacotes
da Microsoft \cite{microsoft_winget}. Ele vem no Windows 11 e nas versões
atuais do Windows 10 (a partir da 1809), como parte do aplicativo
\emph{App Installer}.

\subsection{Conferir o winget}

Abra o \textbf{PowerShell} (menu Iniciar \(\rightarrow\) digite
\enquote{PowerShell}) e rode:

\begin{bashcode}
winget --version
\end{bashcode}

Se aparecer um número de versão (Figura~\ref{fig:winget_installed}), está
pronto. Se o comando não for reconhecido, abra a Microsoft Store, procure
\enquote{App Installer} e atualize ou instale.

\begin{figure}[H]
  \centering
  \includegraphics[width=0.9\textwidth]{./assets/images/07_winget_installed.png}
  \caption{Conferindo o winget}
  \label{fig:winget_installed}
\end{figure}

\begin{dica}[Primeiro uso]
Na primeira vez, o winget pede para aceitar os termos de uso das fontes de
pacotes. Digite \texttt{Y} e aperte \tecla{Enter}.
\end{dica}

\subsection{Instalar o Git}

O \texttt{winget search} mostra o identificador (\emph{Id}) de cada programa
(Figura~\ref{fig:winget_search}). O Id do Git é \texttt{Git.Git}:

\begin{bashcode}
winget search git
winget install --id Git.Git -e --source winget
\end{bashcode}

\begin{figure}[H]
  \centering
  \includegraphics[width=0.9\textwidth]{./assets/images/08_winget_search.png}
  \caption{Buscando o Git no winget}
  \label{fig:winget_search}
\end{figure}

A opção \texttt{-e} exige o Id exato, e \texttt{--source winget} evita a
Microsoft Store. É a forma recomendada pelo site oficial do Git
\cite{git_windows}. Aguarde o fim da instalação
(Figura~\ref{fig:install_git}).

\begin{figure}[H]
  \centering
  \includegraphics[width=0.9\textwidth]{./assets/images/09_install_git.png}
  \caption{Instalação do Git}
  \label{fig:install_git}
\end{figure}

O Git para Windows já traz o \textbf{Git Bash}, o terminal que usaremos daqui
em diante, o \textbf{Git LFS} (capítulo~\ref{cap:lfs}) e o \textbf{Git
Credential Manager}, que guarda o seu login com segurança.

\subsection{Instalar o GitHub CLI}

\begin{bashcode}
winget install --id GitHub.cli
\end{bashcode}

\begin{figure}[H]
  \centering
  \includegraphics[width=0.9\textwidth]{./assets/images/11_install_cli.png}
  \caption{Instalando o GitHub CLI}
  \label{fig:install_cli}
\end{figure}

\subsection{Conferir}

\textbf{Feche o terminal e abra o Git Bash} (menu Iniciar \(\rightarrow\)
\enquote{Git Bash}). O terminal precisa ser reaberto para enxergar os
programas recém-instalados.

\begin{bashcode}
git --version
gh --version
\end{bashcode}

\begin{figure}[H]
  \centering
  \includegraphics[width=0.9\textwidth]{./assets/images/10_git_version.png}
  \caption{Versão do Git}
  \label{fig:git_version}
\end{figure}

\begin{figure}[H]
  \centering
  \includegraphics[width=0.9\textwidth]{./assets/images/12_cli_version.png}
  \caption{Versão do GitHub CLI}
  \label{fig:cli_version}
\end{figure}

\begin{dica}[Atualizar depois]
Para atualizar um programa específico: \texttt{winget upgrade --id Git.Git}.
Evite o \texttt{winget upgrade --all} no meio da aula: ele atualiza todos os
programas do computador e pode demorar bastante.
\end{dica}

% ----------------------------------------------------------------------
\section{Linux (Debian e Ubuntu)}
% ----------------------------------------------------------------------

O Git está nos repositórios da distribuição:

\begin{bashcode}
sudo apt update
sudo apt install git git-lfs
\end{bashcode}

O GitHub CLI tem repositório próprio. Este é o comando oficial
\cite{gh_linux}. Copie inteiro: as barras invertidas no fim das linhas unem
tudo num comando só.

\begin{bashcode}
(type -p wget >/dev/null || (sudo apt update && sudo apt install wget -y)) \
  && sudo mkdir -p -m 755 /etc/apt/keyrings \
  && out=$(mktemp) && wget -nv -O$out https://cli.github.com/packages/githubcli-archive-keyring.gpg \
  && cat $out | sudo tee /etc/apt/keyrings/githubcli-archive-keyring.gpg > /dev/null \
  && sudo chmod go+r /etc/apt/keyrings/githubcli-archive-keyring.gpg \
  && sudo mkdir -p -m 755 /etc/apt/sources.list.d \
  && echo "deb [arch=$(dpkg --print-architecture) signed-by=/etc/apt/keyrings/githubcli-archive-keyring.gpg] https://cli.github.com/packages stable main" | sudo tee /etc/apt/sources.list.d/github-cli.list > /dev/null \
  && sudo apt update \
  && sudo apt install gh -y
\end{bashcode}

\begin{aviso}[Evite o pacote gh da distribuição]
Os mantenedores do GitHub CLI desaconselham o \texttt{gh} empacotado pela
comunidade Debian: algumas versões antigas dele deixaram de funcionar com as
APIs atuais do GitHub \cite{gh_linux}. Use o repositório oficial acima.
\end{aviso}

Outras distribuições (Fedora, Arch, openSUSE) estão em
\url{https://github.com/cli/cli/blob/trunk/docs/install_linux.md}.

% ----------------------------------------------------------------------
\section{macOS}
% ----------------------------------------------------------------------

O caminho mais simples é o \textbf{Homebrew} (\url{https://brew.sh}). Com ele
instalado:

\begin{bashcode}
brew install git git-lfs gh
\end{bashcode}

\begin{dica}
Sem o Homebrew, rodar \texttt{git --version} no Terminal oferece instalar as
\emph{Command Line Tools} da Apple, que trazem o Git. O GitHub CLI, nesse
caso, é instalado pelo pacote em \url{https://cli.github.com}.
\end{dica}

% ----------------------------------------------------------------------
\section{Conferindo nos três sistemas}
% ----------------------------------------------------------------------

\begin{bashcode}
git --version
gh --version
git lfs version
\end{bashcode}

Os três devem responder com um número de versão. Se algum comando não for
reconhecido, feche e reabra o terminal. Se continuar, reinstale a ferramenta.
          % Preparando o ambiente (Windows, Linux, macOS)
% ======================================================================
\chapter{Git e GitHub CLI: configuração e login}
\label{cap:git}
% ======================================================================

\section{O que é o Git}

\begin{conceito}[Git]
O Git é um \textbf{sistema distribuído de controle de versão}, gratuito e de
código aberto. Ele registra cada alteração num \textbf{commit}: uma fotografia
do projeto com autor, data e uma mensagem explicando o porquê.
\enquote{Distribuído} quer dizer que cada pessoa tem o histórico completo no
próprio computador.
\end{conceito}

Um arquivo passa por três lugares:

\begin{table}[H]
  \centering
  \caption{As três áreas do Git}
  \label{tab:tres_areas}
  \begin{tabularx}{\textwidth}{@{}l X l@{}}
    \toprule
    \textbf{Área} & \textbf{O que é} & \textbf{Como chega lá} \\
    \midrule
    Diretório de trabalho & os arquivos como você os vê e edita & editando \\
    Área de preparação (\emph{staging}) & o que vai entrar no próximo commit & \comando{git add} \\
    Repositório & o histórico de commits, na pasta oculta \arq{.git} & \comando{git commit} \\
    \bottomrule
  \end{tabularx}
\end{table}

\section{O que é o GitHub CLI}

O GitHub CLI (\texttt{gh}) faz pelo terminal o que você faria no site: criar
repositórios, abrir e revisar Pull Requests, acompanhar workflows. Ele também
resolve a parte mais chata para iniciantes: o \textbf{login}.

% ----------------------------------------------------------------------
\section{Dizendo ao Git quem é você}
% ----------------------------------------------------------------------

Todo commit leva nome e e-mail. Configure uma vez por computador
(Figura~\ref{fig:git_base_config}):

\begin{bashcode}
git config --global user.name "Seu Nome"
git config --global user.email "seu.email@exemplo.com"
\end{bashcode}

\begin{figure}[H]
  \centering
  \includegraphics[width=0.9\textwidth]{./assets/images/14_git_base_config.png}
  \caption{Configurando nome e e-mail}
  \label{fig:git_base_config}
\end{figure}

\begin{dica}[E-mail cadastrado e privacidade]
Use um e-mail cadastrado na sua conta do GitHub, para os commits aparecerem
ligados ao seu perfil. Se não quiser expor o e-mail pessoal, marque
\menu{Keep my email addresses private} em \menu{Settings} \(\rightarrow\)
\menu{Emails} e use o endereço \texttt{...@users.noreply.github.com} que o
GitHub mostra nessa mesma tela.
\end{dica}

Três ajustes que evitam surpresas:

\begin{bashcode}
# novos repositórios começam na branch main, como no GitHub
git config --global init.defaultBranch main

# git pull faz merge quando as branches divergem
git config --global pull.rebase false

# (opcional) VS Code como editor das mensagens de commit
git config --global core.editor "code --wait"
\end{bashcode}

\begin{aviso}[Por que o pull.rebase false]
Sem essa configuração, um \comando{git pull} com alterações dos dois lados
termina em \texttt{fatal: Need to specify how to reconcile divergent
branches}. O instalador do Git para Windows já define essa opção; no Linux e
no macOS é preciso configurar.
\end{aviso}

Para conferir tudo: \comando{git config --global --list}.

% ----------------------------------------------------------------------
\section{Fazendo login no GitHub}
% ----------------------------------------------------------------------

Desde agosto de 2021, o GitHub não aceita a senha da conta em operações do Git
pela linha de comando. É preciso um token ou uma chave SSH. O jeito mais
simples de configurar isso é o GitHub CLI.

\subsection{Pelo GitHub CLI (recomendado)}

\begin{bashcode}
gh auth login
\end{bashcode}

Responda às perguntas assim
(Figuras~\ref{fig:gh_auth_login_0} a~\ref{fig:gh_auth_login_4}):

\begin{table}[H]
  \centering
  \caption{Respostas do \texttt{gh auth login}}
  \label{tab:gh_auth}
  \begin{tabularx}{\textwidth}{@{}X l@{}}
    \toprule
    \textbf{Pergunta} & \textbf{Resposta} \\
    \midrule
    Where do you use GitHub? & \texttt{GitHub.com} \\
    What is your preferred protocol for Git operations? & \texttt{HTTPS} \\
    Authenticate Git with your GitHub credentials? & \texttt{Yes} \\
    How would you like to authenticate GitHub CLI? & \texttt{Login with a web browser} \\
    \bottomrule
  \end{tabularx}
\end{table}

O terminal mostra um código de uso único. Aperte \tecla{Enter} para abrir o
navegador, entre na conta, digite o código e autorize.

\begin{figure}[H]
  \centering
  \includegraphics[width=0.85\textwidth]{./assets/images/13_0_gh_auth_login.png}
  \caption{\texttt{gh auth login}: escolhendo o GitHub.com}
  \label{fig:gh_auth_login_0}
\end{figure}
\begin{figure}[H]
  \centering
  \includegraphics[width=0.85\textwidth]{./assets/images/13_1_gh_auth_login.png}
  \caption{\texttt{gh auth login}: escolhendo o protocolo}
  \label{fig:gh_auth_login_1}
\end{figure}
\begin{figure}[H]
  \centering
  \includegraphics[width=0.85\textwidth]{./assets/images/13_2_gh_auth_login.png}
  \caption{\texttt{gh auth login}: autenticando o Git com as credenciais do GitHub}
  \label{fig:gh_auth_login_2}
\end{figure}
\begin{figure}[H]
  \centering
  \includegraphics[width=0.85\textwidth]{./assets/images/13_3_gh_auth_login.png}
  \caption{\texttt{gh auth login}: login pelo navegador}
  \label{fig:gh_auth_login_3}
\end{figure}
\begin{figure}[H]
  \centering
  \includegraphics[width=0.85\textwidth]{./assets/images/13_4_gh_auth_login.png}
  \caption{\texttt{gh auth login}: código de uso único}
  \label{fig:gh_auth_login_4}
\end{figure}

\begin{conceito}[O que o \enquote{Yes} faz]
Responder \texttt{Yes} em \enquote{Authenticate Git with your GitHub
credentials?} faz o \texttt{gh} emprestar o login dele ao Git. É isso que faz
o \comando{git push} funcionar sem pedir senha. Se você respondeu
\texttt{No}, rode \comando{gh auth setup-git} para configurar depois
\cite{gh_manual}.
\end{conceito}

Para conferir: \comando{gh auth status}.

\subsection{Por chave SSH (alternativa)}

Com SSH, o computador prova quem é usando um par de chaves: a
\textbf{privada} fica só com você, e a \textbf{pública} é cadastrada no GitHub.

\begin{bashcode}
# cria o par de chaves (aceite o local padrão e escolha uma senha)
ssh-keygen -t ed25519 -C "seu.email@exemplo.com"

# cadastra a chave pública na sua conta
gh ssh-key add ~/.ssh/id_ed25519.pub --title "Meu notebook"

# testa a conexão
ssh -T git@github.com
\end{bashcode}

Na primeira conexão, o SSH pergunta se confia no servidor: responda
\texttt{yes}. A resposta esperada é:

\begin{saida}
Hi seu-usuario! You've successfully authenticated, but GitHub does not provide shell access.
\end{saida}

Com SSH, os endereços dos repositórios passam a ser
\texttt{git@github.com:usuario/repositorio.git}, em vez de
\texttt{https://github.com/...}.

\subsection{Por token de acesso pessoal (para scripts)}

Um \textbf{Personal Access Token} (PAT) é uma senha de uso restrito, com prazo
de validade, que você gera no site. Para uso pessoal no terminal, o
\texttt{gh auth login} já resolve. O PAT faz sentido para scripts e
automações.

\begin{enumerate}
  \item Acesse \url{https://github.com/settings/personal-access-tokens/new}
        (token \textbf{fine-grained}).
  \item Dê um nome, uma validade e escolha \textbf{só os repositórios} e
        \textbf{só as permissões} necessárias.
  \item Gere e copie o valor. Ele não aparece de novo.
\end{enumerate}

\begin{aviso}[Token é senha]
Não compartilhe o token nem o grave em arquivos do projeto. Evite o
\texttt{git config credential.helper store}: ele salva o token em texto puro
no arquivo \arq{\textasciitilde/.git-credentials}. O Git Credential Manager
(Windows) e o \texttt{gh auth setup-git} guardam as credenciais de forma
segura.
\end{aviso}

% ----------------------------------------------------------------------
\section{Comandos do dia a dia}
% ----------------------------------------------------------------------

\begin{table}[H]
  \centering
  \caption{Comandos básicos do Git}
  \label{tab:git_basico}
  \begin{tabularx}{\textwidth}{@{}l X@{}}
    \toprule
    \textbf{Comando} & \textbf{O que faz} \\
    \midrule
    \comando{git init} & transforma a pasta atual num repositório \\
    \comando{git clone <url>} & baixa uma cópia completa de um repositório \\
    \comando{git status} & mostra o que mudou e o que está preparado \\
    \comando{git add <arquivo>} & prepara um arquivo para o próximo commit \\
    \comando{git commit -m "..."} & registra as alterações preparadas \\
    \comando{git log --oneline} & mostra o histórico, um commit por linha \\
    \comando{git switch <branch>} & troca de branch \\
    \comando{git switch -c <branch>} & cria uma branch e já entra nela \\
    \comando{git restore <arquivo>} & descarta alterações não preparadas \\
    \comando{git pull} & baixa e integra as novidades do remoto \\
    \comando{git push} & envia os seus commits para o remoto \\
    \comando{git merge <branch>} & traz as alterações de outra branch \\
    \bottomrule
  \end{tabularx}
\end{table}

\begin{table}[H]
  \centering
  \caption{Comandos básicos do GitHub CLI}
  \label{tab:gh_basico}
  \begin{tabularx}{\textwidth}{@{}l X@{}}
    \toprule
    \textbf{Comando} & \textbf{O que faz} \\
    \midrule
    \comando{gh auth login} & faz o login \\
    \comando{gh auth status} & mostra com qual conta você está logado \\
    \comando{gh repo create} & cria um repositório no GitHub \\
    \comando{gh repo clone <dono>/<repo>} & clona um repositório pelo nome \\
    \comando{gh repo fork <dono>/<repo>} & cria um fork na sua conta \\
    \comando{gh pr create} & abre um Pull Request da branch atual \\
    \comando{gh pr list} & lista os Pull Requests abertos \\
    \comando{gh run list} & mostra as execuções do GitHub Actions \\
    \bottomrule
  \end{tabularx}
\end{table}

As listas completas estão nos Apêndices~\ref{ap:git}
e~\ref{ap:gh}.

\begin{dica}[switch e checkout]
Tutoriais antigos usam \comando{git checkout} para trocar e criar branches.
Ele ainda funciona, mas mistura funções demais. O \comando{git switch} (para
branches) e o \comando{git restore} (para arquivos) foram criados para
separar essas tarefas e são mais seguros para quem está começando.
\end{dica}
          % Git e GitHub CLI: configuração e autenticação
% ======================================================================
\chapter{Git LFS: arquivos grandes}
\label{cap:lfs}
% ======================================================================

\section{O problema}

O Git foi feito para texto, como código-fonte. Com arquivos grandes e binários
(vídeos, imagens pesadas, modelos de aprendizado de máquina, arquivos
\texttt{.psd} ou \texttt{.zip}), aparecem três problemas:

\begin{itemize}
  \item \textbf{Limite do GitHub:} o GitHub avisa a partir de 50 MiB e
        \textbf{recusa arquivos acima de 100 MiB} num push comum
        \cite{github_lfs}. Esse limite é do GitHub, não do Git.
  \item \textbf{Histórico que só cresce:} binários quase não se beneficiam da
        forma como o Git guarda diferenças entre versões. Cada versão nova de
        um arquivo grande ocupa quase o tamanho inteiro de novo, e todo
        \comando{git clone} baixa todas elas.
  \item \textbf{Operações lentas:} com muitos binários no histórico,
        \comando{clone}, \comando{fetch} e \comando{checkout} ficam lentos.
\end{itemize}

\section{A solução}

\begin{conceito}[Git LFS]
O \textbf{Git LFS} (\emph{Large File Storage}) é uma extensão do Git. No
repositório fica só um \textbf{ponteiro} pequeno, de texto. O conteúdo real vai
para um armazenamento separado, no próprio GitHub, e é baixado só quando
necessário. Os comandos do dia a dia (\comando{add}, \comando{commit},
\comando{push}) continuam os mesmos.
\end{conceito}

O ponteiro gravado no Git tem esta cara:

\begin{saida}
version https://git-lfs.github.com/spec/v1
oid sha256:e1ea497dccb8bd9d842c37f70699210c4db6d6981d7c41fecd1d028d16fade26
size 200000
\end{saida}

\section{Instalação}

\begin{itemize}
  \item \textbf{Windows:} já vem com o Git para Windows.
  \item \textbf{Linux (Debian e Ubuntu):} \comando{sudo apt install git-lfs}.
  \item \textbf{macOS:} \comando{brew install git-lfs}.
\end{itemize}

Depois, rode \textbf{uma vez por usuário}:

\begin{bashcode}
git lfs install
\end{bashcode}

\begin{saida}
Updated git hooks.
Git LFS initialized.
\end{saida}

Esse comando não instala o programa: ele liga o LFS na configuração do seu
Git.

\section{Usando}

\begin{bashcode}
# diz ao LFS quais arquivos ele deve cuidar
git lfs track "*.zip"
git lfs track "*.psd"

# a regra fica gravada no .gitattributes, que também vai para o repositório
git add .gitattributes
git add materiais.zip
git commit -m "chore: versiona os materiais com Git LFS"
git push

# confere o que está no LFS
git lfs ls-files
\end{bashcode}

\begin{aviso}[Configure antes de adicionar]
O \comando{git lfs track} só vale para o que for adicionado \emph{depois}. Um
arquivo grande que já entrou no histórico como arquivo comum continua lá.
Mover arquivos antigos para o LFS exige \comando{git lfs migrate}, que
\textbf{reescreve o histórico} e precisa ser combinado com a equipe.
\end{aviso}

\section{Limites e custos no GitHub}

\begin{itemize}
  \item Cada conta tem \textbf{10 GiB de armazenamento} e \textbf{10 GiB de
        banda} de LFS por mês, tanto no plano Free quanto no Pro
        \cite{github_lfs_uso}. Acima disso, o uso é cobrado.
  \item O tamanho máximo de um arquivo no LFS é de 2 GB nos planos Free e
        Pro.
  \item \textbf{Apagar um arquivo do LFS não libera espaço} sozinho. Para
        remover de verdade, é preciso reescrever o histórico ou apagar o
        repositório.
\end{itemize}

\begin{dica}[Quando não usar]
Não use o LFS para arquivos pequenos em grande quantidade nem para o que pode
ser gerado de novo (arquivos compilados, dependências): esses devem ir para o
\arq{.gitignore}.
\end{dica}
          % Git LFS
% ======================================================================
\chapter{Commits, merges e Pull Requests}
\label{cap:commits}
% ======================================================================

\section{Introdução}
Commits, merges e pull requests são elementos centrais no fluxo de trabalho com Git e plataformas como GitHub.  
Um \textbf{commit} é uma unidade lógica de alteração no repositório; um \textbf{merge} integra mudanças de uma branch em outra; e um \textbf{pull request} (PR) é uma solicitação formal de revisão e integração de uma branch — normalmente usada em colaboração para revisar código, executar checks automáticos (CI) e registrar a decisão de integrar.

\section{Commits}

\subsection{O que é um commit}
Um commit registra o estado do diretório de trabalho (os arquivos staged) em um nó do histórico do Git. Cada commit tem um identificador (SHA), metadata (autor, data) e uma mensagem que descreve a mudança.

\subsection{Boas práticas de commits}
\begin{itemize}
  \item Faça commits atômicos: cada commit deve representar uma \emph{única} mudança lógica.  
  \item Mensagens claras: use assunto imperativo curto (\(\sim\)50 caracteres) + linha em branco + corpo explicativo se necessário (72 caracteres por linha).  
  \item Inclua referência a issues quando relevante (ex.: ``Closes \#42'').  
  \item Evite commitar arquivos gerados (binários, dependências) — use \texttt{.gitignore}.  
  \item Assine os commits (capítulo~\ref{cap:assinatura}).
\end{itemize}

\subsection{Fazendo commits — passo a passo}
\begin{enumerate}
  \item Verificar o estado dos arquivos:
  \begin{bashcode}
  git status
  \end{bashcode}
  \item Ver as mudanças não staged:
  \begin{bashcode}
  git diff       # diferenças não adicionadas ao stage
  git diff --staged  # diferenças preparadas para commit
  \end{bashcode}
  \item Adicionar alterações ao stage (todo arquivo ou interativo):
  \begin{bashcode}
  git add caminho/arquivo.txt  # adiciona arquivo específico
  git add .                    # adiciona tudo (cuidado)
  git add -p                   # adiciona parcialmente (patch)
  \end{bashcode}
  \item Criar o commit com boa mensagem:
  \begin{bashcode}
  git commit -m "Assunto curto em imperativo"
  # ou para mensagem longa (editor):
  git commit
  \end{bashcode}
  Exemplo de mensagem:
  \begin{bashcode}
  Corrige cálculo de juros

  Ajusta a fórmula de cálculo para considerar juros compostos
  quando o período é maior que 12 meses. Testes unitários
  adicionados para cobrir casos de fronteira.
  \end{bashcode}
  \item Ver histórico resumido:
  \begin{bashcode}
  git log --oneline --graph --decorate --all
  \end{bashcode}
\end{enumerate}

\subsection{Editar o último commit / corrigir mensagens}
\begin{bashcode}
# alterar o conteúdo do último commit (já staged)
git commit --amend

# alterar apenas a mensagem do último commit
git commit --amend -m "Nova mensagem"
\end{bashcode}
\textbf{Atenção:} se o commit já foi enviado ao remoto, evite \texttt{--amend} sem combinar com a equipe — ele reescreve o histórico e exigirá push forçado.

\subsection{Desfazer / alterar staging}
\begin{bashcode}
git restore --staged arquivo.txt   # remove do stage, mantém
                                   # alteração no working tree

git restore arquivo.txt            # descarta alteração no
                                   # working tree (se não commitada)

git reset --soft HEAD~1            # desfaz último commit, mantendo
                                   # mudanças staged

git reset --mixed HEAD~1           # desfaz último commit, mantendo
                                   # mudanças no working tree
                                   # (unstaged)

git reset --hard HEAD~1            # desfaz último commit e 
                                   # descarta mudanças (CUIDADO)
\end{bashcode}

\begin{dica}[Como escrever a mensagem]
O capítulo~\ref{cap:padroes} mostra um padrão para as mensagens de commit, o \emph{Conventional Commits}.
\end{dica}

\section{Merges}

\subsection{Tipos de merge}
\begin{description}
  \item[Fast-forward] Se a branch destino não avançou desde que a feature foi criada, o Git apenas avança o ponteiro (sem novo commit de merge).
  \item[Merge commit] Cria um commit de merge que documenta a união de dois históricos (útil para preservar contexto de branch).
  \item[Rebase] Reaplica commits de uma branch sobre outra, produzindo um histórico linear (reduz ``ruído'' dos merges, mas reescreve histórico).
\end{description}

\subsection{Merge local com merge commit (passo a passo)}
\begin{enumerate}
  \item Atualize a branch principal:
  \begin{bashcode}
  git switch main
  git pull origin main
  \end{bashcode}
  \item Mudar para a branch de feature (se necessário):
  \begin{bashcode}
  git switch feature/minha-feature
  git pull origin feature/minha-feature
  \end{bashcode}
  \item Voltar para a branch de destino e fazer merge:
  \begin{bashcode}
  git switch main
  git merge --no-ff feature/minha-feature
  # --no-ff força um commit de merge,
  # preservando histórico da branch
  \end{bashcode}
  \item Caso não haja conflitos, o Git criará o commit de merge automaticamente. Em caso de conflitos, siga a seção "Resolver conflitos" abaixo.
  \item Envie as mudanças para o remoto:
  \begin{bashcode}
  git push origin main
  \end{bashcode}
\end{enumerate}

\subsection{Rebase (passo a passo) — para um histórico linear}
\begin{enumerate}
  \item Atualize a base:
  \begin{bashcode}
  git switch main
  git pull origin main
  \end{bashcode}
  \item Rebase da feature sobre a main:
  \begin{bashcode}
  git switch feature/minha-feature
  git rebase main
  \end{bashcode}
  \item Resolva conflitos (se aparecerem), use \texttt{git rebase --continue} e, ao final:
  \begin{bashcode}
  # force push (com cuidado) após reescrever histórico
  git push --force-with-lease origin feature/minha-feature
  \end{bashcode}
  \item \textbf{Observação:} reescrever histórico (rebase + push forçado) requer coordenação com outros colaboradores.
\end{enumerate}

\subsection{Resolver conflitos — passo a passo}
\begin{enumerate}
  \item Ao encontrar conflito durante merge/rebase, o Git mostra arquivos conflitantes:
  \begin{bashcode}
  CONFLICT (content): Merge conflict in caminho/arquivo.txt
  \end{bashcode}
  \item Abra o arquivo e localize os marcadores:
  \begin{bashcode}
  <<<<<<< HEAD
  conteúdo na branch atual (main)
  =======
  conteúdo vindo da branch feature/minha-feature
  >>>>>>> feature/minha-feature
  \end{bashcode}
  \item Edite o arquivo para a versão desejada (manter, combinar ou reescrever o trecho).
  \item Marque o conflito como resolvido:
  \begin{bashcode}
  git add caminho/arquivo.txt
  # se for rebase:
  git rebase --continue
  # se for merge:
  git commit   # caso o Git não tenha criado o commit automaticamente
  \end{bashcode}
  \item Se quiser abortar a operação:
  \begin{bashcode}
  git merge --abort    # durante um merge
  git rebase --abort   # durante um rebase
  \end{bashcode}
\end{enumerate}

\subsection{Squash e reescrita de commits (passo a passo)}
\begin{enumerate}
  \item Rebase interativo para combinar commits locais:
  \begin{bashcode}
  git switch feature/minha-feature
  git rebase -i main
  \end{bashcode}
  \item No editor que abre, marque \texttt{squash} (ou \texttt{s}) nos commits que deseja unir ao commit anterior. Salve e feche.
  \item Após o rebase, force push com segurança:
  \begin{bashcode}
  git push --force-with-lease origin feature/minha-feature
  \end{bashcode}
  \item Use \texttt{--force-with-lease} em vez de \texttt{--force} quando possível — ele evita sobrescrever pushes alheios.
\end{enumerate}

\section{Pull Requests (PR)}

\subsection{O que é um Pull Request}
Um PR é uma solicitação para integrar as mudanças de uma branch em outra (normalmente de uma branch de feature para \texttt{main} ou \texttt{develop}) e serve como ponto central para revisão de código, execução de pipelines de CI e documentação da mudança.

\subsection{Fluxo básico — criando um PR (via web)}
\begin{enumerate}
  \item Crie uma branch local para a sua feature:
  \begin{bashcode}
  git switch -c feature/minha-feature
  \end{bashcode}
  \item Faça commits locais e envie a branch para o remoto:
  \begin{bashcode}
  git push -u origin feature/minha-feature
  \end{bashcode}
  \item No repositório do GitHub, acesse \texttt{Pull requests} \(\rightarrow\) \texttt{New pull request}.
  \item Selecione a \textbf{branch de origem} (feature/minha-feature) e a \textbf{branch de destino} (ex.: main).
  \item Preencha o título e a descrição: explique o \textbf{porquê} e o \textbf{o que} foi alterado. Use referências a issues (ex.: ``Closes \#123'').
  \item Configure revisores, labels, milestone e assignees.
  \item Se ainda não está pronta, crie como \textbf{Draft pull request} (rascunho).
  \item Aguarde revisão, corrija comentários fazendo novos commits na mesma branch e push — eles serão anexados automaticamente ao PR.
  \item Quando aprovado, escolha a estratégia de merge (merge commit / squash and merge / rebase and merge) e realize o merge.
\end{enumerate}

\subsection{Criar e gerenciar PRs via GitHub CLI (passo a passo)}
\begin{bashcode}
# autentique (uma vez)
gh auth login

# depois de push da branch
gh pr create --base main --head feature/minha-feature --title "Título do PR" --body "Descrição detalhada" --reviewer usuario1,usuario2

# abrir PR como draft:
gh pr create --draft --base main --head feature/minha-feature --fill

# listar PRs locais:
gh pr list

# fechar ou mesclar via CLI:
gh pr merge <numero-ou-url> --squash --delete-branch
# ou
gh pr merge <numero> --merge      # cria commit de merge
gh pr merge <numero> --rebase     # rebase and merge
\end{bashcode}

\subsection{Checklist para revisão de Pull Request}
\begin{itemize}
  \item O PR tem um título e descrição claros (o \emph{porquê} e o \emph{o que});  
  \item Testes automatizados adicionados/atualizados e pipeline CI passando;  
  \item Código segue padrões de lint e estilo;  
  \item Mudanças pequenas e focadas (um PR por responsabilidade);  
  \item Documentação e comentários quando necessário;  
  \item Evidências visuais (screenshots) quando há alterações de UI.  
\end{itemize}

\subsection{Depois do merge — limpeza e sincronização}
\begin{bashcode}
# atualizar a branch principal local
git switch main
git pull origin main

# remover branch remota (após merge)
git push origin --delete feature/minha-feature

# remover branch local
git branch -d feature/minha-feature

# caso a branch não possa ser deletada por não estar totalmente mesclada:
git branch -D feature/minha-feature
\end{bashcode}
Também é útil rodar:
\begin{bashcode}
git fetch --prune
\end{bashcode}
para remover referências remotas deletadas.

\section{Boas práticas e recomendações finais}
\begin{itemize}
  \item Mantenha PRs pequenos e revisáveis; grandes PRs demoram mais para receber feedback.  
  \item Automatize checks com CI (testes, lint, análise estática) e impeça merge enquanto falharem (branch protection).  
  \item Use \texttt{--force-with-lease} quando precisar reescrever histórico; nunca force sem checar se colegas não empurraram commits.  
  \item Documente o fluxo do time (merge strategy preferida — ex.: squash para commits limpos, merge commit para histórico preservado).  
  \item Escreva mensagens de commit úteis e mantenha um padrão no time (ex.: Conventional Commits) para facilitar geração automática de changelogs.
\end{itemize}

\section{Exemplos rápidos de comandos úteis}
\begin{bashcode}
# ver status e diferenças
git status
git diff
git diff --staged

# histórico e log
git log --oneline --graph --decorate --all

# adicionar parcialmente
git add -p

# rebase interativo (últimos 3 commits)
git rebase -i HEAD~3

# desfazer mudanças locais (trabalhe com cuidado)
git restore arquivo.txt
git reset --hard HEAD

# forçar push com segurança
git push --force-with-lease origin minha-branch
\end{bashcode}

% Fim do capítulo          % Commits, merges e Pull Requests
% ======================================================================
\chapter{Padrões de commit}
\label{cap:padroes}

No capítulo~\ref{cap:commits} vimos como fazer commits. Este capítulo trata de \textbf{como escrever essa mensagem} de um jeito que a equipe inteira (e as ferramentas) consigam ler.

\section{Por que padronizar}

Compare dois históricos do mesmo projeto:

\begin{saida}
a41aded ajustes
c2524a8 arrumei
afe16bb agora vai
4203776 mudanças do joão
\end{saida}
\begin{saida}
a41aded fix(login): corrige senha vazia
c2524a8 feat(relatorio): exporta em PDF
afe16bb docs: explica a instalação
4203776 test(login): cobre senha vazia
\end{saida}

No segundo, dá para saber \textbf{o tipo} da mudança, \textbf{onde} ela aconteceu e \textbf{o que} foi feito, sem abrir um arquivo sequer. Um padrão de commit também permite automatizar tarefas: gerar notas de versão, decidir o próximo número de versão e barrar mensagens fora do combinado.

\begin{conceito}[A regra mais importante]
O melhor padrão é \textbf{o que a equipe combinou e segue}. Os padrões deste capítulo são os mais usados, mas consistência vale mais que qualquer escolha isolada.
\end{conceito}

\section{Conventional Commits}

O padrão mais difundido é o \textbf{Conventional Commits} 1.0.0 (\url{https://www.conventionalcommits.org/pt-br/v1.0.0/}). A estrutura é:

\begin{bashcode}
<tipo>[escopo opcional]: <descrição>

[corpo opcional]

[rodapé(s) opcional(is)]
\end{bashcode}

As regras principais da especificação:

\begin{itemize}
  \item o \textbf{tipo} é obrigatório e vem primeiro, seguido de \textbf{dois-pontos e espaço};
  \item o \textbf{escopo} é opcional: um substantivo entre parênteses que diz qual parte do sistema mudou, como \arq{fix(parser):};
  \item a \textbf{descrição} vem logo depois de \arq{: } e resume a mudança;
  \item o \textbf{corpo} é opcional e começa \textbf{uma linha em branco} depois da descrição, com quantos parágrafos forem necessários;
  \item os \textbf{rodapés} são opcionais e vêm uma linha em branco depois do corpo, no formato \arq{Token: valor} (por exemplo, \arq{Refs: \#123} ou \arq{Reviewed-by: Maria});
  \item uma \textbf{quebra de compatibilidade} (\emph{breaking change}) é indicada com \arq{!} antes dos dois-pontos, com um rodapé \arq{BREAKING CHANGE: <explicação>}, ou com os dois.
\end{itemize}

\begin{bashcode}
feat(api)!: exige autenticação em todos os endpoints

Antes, os endpoints de consulta eram públicos. Agora todo pedido
precisa do cabeçalho Authorization com um token válido.

BREAKING CHANGE: clientes sem token passam a receber 401.
Refs: #42
\end{bashcode}

\subsection{Os tipos}

A especificação só define dois tipos: \textbf{\arq{feat}} (novo recurso) e \textbf{\arq{fix}} (correção de bug). Os outros são convenção, herdada das regras do projeto Angular. A lista abaixo é a aceita pela configuração padrão do \emph{commitlint} (seção \ref{sec:commitlint}):

\begin{center}
\renewcommand{\arraystretch}{1.3}
\begin{tabularx}{\linewidth}{@{}>{\ttfamily\color{roxoescuro}}l X@{}}
\toprule
\textrm{\textbf{Tipo}} & \textbf{Quando usar} \\
\midrule
feat & um recurso novo para quem usa o sistema \\
fix & a correção de um bug \\
docs & só documentação (README, comentários, manuais) \\
style & formatação que não muda o comportamento (espaços, ponto e vírgula, indentação) \\
refactor & reorganização do código sem mudar o que ele faz nem corrigir bug \\
perf & melhoria de desempenho \\
test & criação ou ajuste de testes \\
build & sistema de build ou dependências (npm, pip, Maven\ldots) \\
ci & configuração de integração contínua (GitHub Actions) \\
chore & tarefas de manutenção que não se encaixam nos outros tipos \\
revert & desfaz um commit anterior \\
\bottomrule
\end{tabularx}
\end{center}

\begin{dica}[Tipos e versões]
Com Conventional Commits, ferramentas conseguem calcular a próxima versão pelo \emph{Versionamento Semântico}: \arq{fix} sobe o número de correção (1.2.\textbf{3}), \arq{feat} sobe o número de recurso (1.\textbf{3}.0) e um \arq{BREAKING CHANGE} sobe o número principal (\textbf{2}.0.0).
\end{dica}

\begin{aviso}[Tipos fora da lista]
Algumas referências populares em português acrescentam tipos como \arq{raw}, \arq{cleanup}, \arq{remove} e \arq{init}. Eles não fazem parte da especificação nem da configuração padrão do commitlint, que os rejeita. Só use se a equipe combinar e configurar isso.
\end{aviso}

\section{Boas práticas para a mensagem}

\begin{itemize}
  \item \textbf{Uma mudança por commit.} Se a descrição precisa de um \enquote{e}, provavelmente são dois commits. Commits pequenos são mais fáceis de revisar e de desfazer.
  \item \textbf{Descrição curta.} O livro \emph{Pro Git} recomenda uma primeira linha de até uns 50 caracteres e o corpo quebrado em 72 colunas. A configuração padrão do commitlint aceita até 100 caracteres no cabeçalho.
  \item \textbf{O corpo explica o porquê.} O \enquote{o quê} está no \comando{git diff}. O motivo da mudança só existe na mensagem.
  \item \textbf{Um único modo verbal.} As convenções em inglês usam o imperativo (\enquote{add}, \enquote{fix}). Em português, equipes usam o imperativo (\enquote{corrige}), o presente (\enquote{corrige}) ou o gerúndio (\enquote{corrigindo}). A especificação não exige nenhum. Escolha um e mantenha.
  \item \textbf{Idioma combinado.} A especificação também não fixa idioma. Misturar português e inglês no mesmo histórico atrapalha a busca.
\end{itemize}

\section{Gitmoji: emojis nas mensagens}

O \textbf{gitmoji} (\url{https://gitmoji.dev}) associa um emoji a cada tipo de mudança. No texto do commit, o emoji aparece como um código entre dois-pontos, o \emph{shortcode}. Alguns dos mais usados, com a descrição oficial traduzida:

\begin{center}
\renewcommand{\arraystretch}{1.25}
\small
\begin{tabularx}{\linewidth}{@{}>{\ttfamily\color{roxoescuro}}l X >{\ttfamily\color{roxoescuro}}l X@{}}
\toprule
\textrm{\textbf{Código}} & \textbf{Significado} & \textrm{\textbf{Código}} & \textbf{Significado} \\
\midrule
:sparkles: & novo recurso & :bug: & corrige um bug \\
:memo: & documentação & :art: & estrutura ou formato do código \\
:zap: & desempenho & :recycle: & refatoração \\
:fire: & remove código ou arquivos & :truck: & move ou renomeia arquivos \\
:white\_check\_mark: & testes & :construction\_worker: & CI \\
:wrench: & arquivos de configuração & :package: & arquivos compilados ou pacotes \\
:lock: & segurança ou privacidade & :ambulance: & correção crítica urgente \\
:pencil2: & corrige erro de digitação & :rewind: & reverte mudanças \\
:tada: & começa um projeto & :bookmark: & versão ou tag \\
:page\_facing\_up: & licença & :card\_file\_box: & mudanças de banco de dados \\
\bottomrule
\end{tabularx}
\end{center}

A lista completa tem 75 emojis e está em \url{https://gitmoji.dev}.

\begin{aviso}[Três cuidados com emojis]
\begin{enumerate}
  \item \textbf{O terminal não converte.} O GitHub mostra o emoji na página do commit, mas o \comando{git log} exibe o texto \arq{:sparkles:} como foi escrito.
  \item \textbf{O mesmo emoji muda de sentido entre fontes.} Na lista oficial, \arq{:boom:} é quebra de compatibilidade, mas há tutoriais que o usam para \enquote{reverter}. Siga a lista oficial.
  \item \textbf{Emoji antes do tipo quebra o padrão.} A especificação exige que a mensagem comece pelo tipo. Uma mensagem como \arq{:bug: fix: corrige o login} é rejeitada pelo commitlint na configuração padrão.
\end{enumerate}
\end{aviso}

\section{Um padrão combinando os dois}

Muitas equipes juntam gitmoji e Conventional Commits num padrão próprio. É o caso dos repositórios do autor deste guia, que usam:

\begin{bashcode}
<:emoji:> <tipo>(<escopo>): <Descrição no gerúndio>

:memo: docs(guia): Adicionando os exercícios de assinatura
:bug: fix(login): Corrigindo o acesso com senha vazia
:truck: chore(docs): Movendo o manual para a pasta materials
\end{bashcode}

Não há nada de errado nisso, desde que \textbf{seja uma decisão da equipe, esteja escrita} (no README ou num guia de contribuição) e \textbf{as ferramentas estejam configuradas para aceitá-la}. O ponto é a consistência.

\section{Ferramentas}
\label{sec:commitlint}

\subsection*{commitlint: barrar mensagens fora do padrão}

O \textbf{commitlint} (\url{https://commitlint.js.org}) confere cada mensagem antes de o commit ser gravado. Num projeto com Node.js:

\begin{bashcode}
npm install -D @commitlint/cli @commitlint/config-conventional
echo "export default { extends: ['@commitlint/config-conventional'] };" > commitlint.config.js

# testar uma mensagem sem precisar fazer commit
echo "feat: adiciona exportação em PDF" | npx commitlint
\end{bashcode}

Para rodar automaticamente em todo \comando{git commit}, o guia oficial usa o \emph{husky}:

\begin{bashcode}
npm install --save-dev husky
npx husky init
echo "npx --no -- commitlint --edit \$1" > .husky/commit-msg
\end{bashcode}

Depois disso, \comando{git commit -m "foo: teste"} é recusado, porque \arq{foo} não é um tipo válido.

\begin{aviso}[A configuração padrão é rígida]
Além de exigir o tipo no início, a configuração padrão rejeita descrições que começam com letra maiúscula. Um padrão como o da seção anterior precisa de regras próprias no \arq{commitlint.config.js}.
\end{aviso}

\subsection*{gitmoji-cli: escolher o emoji num menu}

\begin{bashcode}
npm i -g gitmoji-cli
gitmoji -c        # monta o commit de forma interativa
gitmoji -l        # lista os emojis
\end{bashcode}

\subsection*{Sem ferramenta nenhuma}

Nada disso é obrigatório. Um arquivo \arq{CONTRIBUTING.md} com o padrão e três exemplos já resolve boa parte do problema, e a revisão do Pull Request (capítulo~\ref{cap:commits}) é um bom momento para lembrar o combinado.
          % Padrões de commit
% ======================================================================
\chapter{GitHub Pages e GitHub Actions}
\label{cap:pages}
% ======================================================================

\section{GitHub Pages: um site a partir do repositório}

O \textbf{GitHub Pages} hospeda \textbf{sites estáticos} (HTML, CSS e
JavaScript) direto de um repositório, de graça, no endereço
\texttt{https://<seu-usuario>.github.io/<repositorio>/}
\cite{github_pages}.

\begin{itemize}
  \item Em repositório \textbf{público}, funciona em qualquer conta. Em
        repositório \textbf{privado}, exige o GitHub Pro (que vem com os
        benefícios de estudante, capítulo~\ref{cap:github}).
  \item Aceita domínio próprio, além do \texttt{github.io}.
  \item Não roda código no servidor: nada de PHP, banco de dados ou Python do
        lado do servidor.
\end{itemize}

\subsection{As duas formas de publicar}

A escolha fica em \menu{Settings} \(\rightarrow\) \menu{Pages}
\(\rightarrow\) \menu{Build and deployment} \(\rightarrow\) \menu{Source}:

\begin{table}[H]
  \centering
  \caption{Fontes de publicação do GitHub Pages}
  \label{tab:pages_fontes}
  \begin{tabularx}{\textwidth}{@{}l X@{}}
    \toprule
    \textbf{Source} & \textbf{Quando usar} \\
    \midrule
    \menu{Deploy from a branch} & o site já está pronto no repositório: basta
      escolher a branch (por exemplo, \texttt{main}) e a pasta
      (\texttt{/(root)} ou \texttt{/docs}) \\
    \menu{GitHub Actions} & o site precisa ser gerado antes de publicar, ou
      você quer controlar o que é publicado. É o caminho recomendado pela
      documentação atual \\
    \bottomrule
  \end{tabularx}
\end{table}

\begin{aviso}[O index.html precisa estar na pasta publicada]
O Pages abre o \arq{index.html} da pasta escolhida. Se ele estiver dentro de
uma subpasta (por exemplo, \arq{jogo-da-memoria/index.html}), o endereço
principal responde com erro 404.
\end{aviso}

% ----------------------------------------------------------------------
\section{GitHub Actions: automação dentro do GitHub}
% ----------------------------------------------------------------------

O \textbf{GitHub Actions} executa tarefas automaticamente quando algo acontece
no repositório: rodar testes a cada push, compilar um projeto, publicar um
site \cite{github_actions}.

\begin{table}[H]
  \centering
  \caption{Vocabulário do GitHub Actions}
  \label{tab:actions_termos}
  \begin{tabularx}{\textwidth}{@{}l X@{}}
    \toprule
    \textbf{Termo} & \textbf{O que é} \\
    \midrule
    Workflow & o processo automatizado, descrito num arquivo YAML em
      \arq{.github/workflows/} \\
    Evento (\texttt{on}) & o que dispara o workflow: um \texttt{push}, um Pull
      Request, um horário, um clique manual (\texttt{workflow\_dispatch}) \\
    Job & um conjunto de passos que roda numa mesma máquina \\
    Step & um passo: um comando (\texttt{run}) ou uma ação pronta
      (\texttt{uses}) \\
    Runner & a máquina virtual do GitHub que executa o job, como
      \texttt{ubuntu-latest} \\
    Action & um passo pronto e reutilizável, como
      \texttt{actions/checkout} \\
    \bottomrule
  \end{tabularx}
\end{table}

\begin{dica}[Minutos]
Em repositórios \textbf{públicos}, o GitHub Actions é gratuito. Em privados,
os minutos por mês dependem do plano (tabela~\ref{tab:free_pro}).
\end{dica}

\subsection{Exemplo 1: publicar um site pelo Actions}

Salve como \arq{.github/workflows/pages.yml} e, em \menu{Settings}
\(\rightarrow\) \menu{Pages}, escolha a Source \menu{GitHub Actions}:

\begin{yamlcode}
name: Publicar o site

on:
  push:
    branches: [ main ]
  workflow_dispatch:

permissions:
  contents: read
  pages: write
  id-token: write

concurrency:
  group: pages
  cancel-in-progress: false

jobs:
  publicar:
    runs-on: ubuntu-latest
    environment:
      name: github-pages
      url: ${{ steps.deploy.outputs.page_url }}
    steps:
      - uses: actions/checkout@v7
      - uses: actions/configure-pages@v6
      - uses: actions/upload-pages-artifact@v5
        with:
          path: .
      - id: deploy
        uses: actions/deploy-pages@v5
\end{yamlcode}

O \texttt{path: .} publica a raiz do repositório. Se o site estiver numa
subpasta, troque pelo caminho dela.

\subsection{Exemplo 2: rodar um script e guardar o resultado}

Este workflow roda o \arq{main.py} (disponível em \arq{materials/}), que cria
um arquivo \arq{resultado.txt}, e grava esse arquivo de volta no repositório:

\begin{yamlcode}
name: Executar o script Python

on:
  push:
    branches: [ main ]
  workflow_dispatch:

permissions:
  contents: write   # necessário para o git push no último passo

jobs:
  executar:
    runs-on: ubuntu-latest
    steps:
      # 1. Baixa o repositório para a máquina do GitHub
      - uses: actions/checkout@v7

      # 2. Instala o Python
      - uses: actions/setup-python@v7
        with:
          python-version: '3.x'

      # 3. Roda o script, que cria o resultado.txt
      - name: Executar o main.py
        run: python main.py

      # 4. Grava o resultado no repositório
      - name: Commitar o resultado
        run: |
          git config user.name "github-actions[bot]"
          git config user.email "github-actions[bot]@users.noreply.github.com"
          git add resultado.txt
          git commit -m "chore: atualiza o resultado.txt" || echo "Nada para commitar"
          git push
\end{yamlcode}

\begin{aviso}[Sem permissão, o push falha]
Por padrão, o token que o workflow recebe (\texttt{GITHUB\_TOKEN}) pode ser
só de leitura. Sem o \texttt{permissions: contents: write}, o
\texttt{git push} do último passo falha com erro 403.
\end{aviso}

\subsection{E o \texttt{[skip ci]}?}

Poderia parecer que o push feito pelo próprio workflow dispara o workflow de
novo, num laço infinito. \textbf{Não dispara}: pushes feitos com o
\texttt{GITHUB\_TOKEN} não iniciam novos workflows. O \texttt{[skip ci]} na
mensagem do commit (ou \texttt{[ci skip]}, \texttt{[no ci]},
\texttt{[skip actions]}) só é necessário quando o push é feito com um token
pessoal ou de aplicativo \cite{github_skip}.

\subsection{Acompanhando}

A aba \menu{Actions} do repositório mostra cada execução, com o log de cada
passo. Pelo terminal:

\begin{bashcode}
gh run list          # últimas execuções
gh run watch         # acompanha a execução em andamento
\end{bashcode}
          % GitHub Pages e GitHub Actions
% ======================================================================
\chapter{Assinatura de commits com GPG e SSH}
\label{cap:assinatura}

\section{O problema: o autor de um commit é só um texto}

No capítulo~\ref{cap:git} configuramos o nome e o e-mail com \comando{git config}. O Git aceita \textbf{qualquer} valor ali, sem conferir nada. Veja:

\begin{bashcode}
git -c user.name="Linus Torvalds" -c user.email="torvalds@linux-foundation.org" \
    commit --allow-empty -m "Commit que parece do Linus"
git log --format='%an <%ae>' -1
\end{bashcode}
\begin{saida}
Linus Torvalds <torvalds@linux-foundation.org>
\end{saida}

O commit foi feito por nós, mas o histórico diz que foi o criador do Git. Se esse commit for enviado ao GitHub, ele aparece com o nome e a foto de quem tem aquele e-mail cadastrado.

\begin{conceito}[Assinatura de commit]
Assinar um commit é anexar a ele uma \textbf{assinatura criptográfica} feita com uma chave privada que só você tem. Quem tem a sua chave \textbf{pública} consegue conferir que o commit veio de você e que não foi alterado depois. O GitHub faz essa conferência e mostra um selo ao lado do commit.
\end{conceito}

\subsection*{Os selos}

\begin{center}
\renewcommand{\arraystretch}{1.35}
\begin{tabularx}{\linewidth}{@{}>{\bfseries}l X@{}}
\toprule
\textbf{Selo} & \textbf{Significado} \\
\midrule
\color{verdeok}Verified & o commit está assinado e a assinatura foi conferida com uma chave cadastrada na conta \\
\color{vermelhoerro}Unverified & o commit está assinado, mas a assinatura não pôde ser conferida (chave não cadastrada, e-mail diferente, chave de outra pessoa) \\
\textrm{(sem selo)} & o commit não foi assinado \\
\bottomrule
\end{tabularx}
\end{center}

\begin{aviso}[O e-mail precisa bater]
Para o selo \textbf{Verified}, o e-mail do commit (\comando{git config user.email}) precisa ser um e-mail \textbf{verificado} na sua conta do GitHub. No caso do GPG, ele também precisa estar dentro da chave. Essa é a causa mais comum de \enquote{Unverified}.
\end{aviso}

\section{GPG ou SSH?}

Há dois jeitos de assinar, e os dois são aceitos pelo GitHub:

\begin{center}
\renewcommand{\arraystretch}{1.35}
\small
\begin{tabularx}{\linewidth}{@{}>{\bfseries\color{roxoescuro}}l >{\raggedright\arraybackslash}X >{\raggedright\arraybackslash}X@{}}
\toprule
 & \textbf{GPG} & \textbf{SSH} \\
\midrule
Ferramenta & GnuPG (\comando{gpg}) & OpenSSH (\comando{ssh-keygen}), que já vem com o Git \\
Requisito & GnuPG instalado & Git 2.34 ou mais novo \\
Chave & uma chave própria, que carrega nome e e-mail & a mesma chave SSH do capítulo~\ref{cap:git} pode servir \\
Validade & pode ter data de expiração embutida & sem data de expiração \\
Conferir no terminal & direto, com o chaveiro do GPG & precisa de um arquivo \arq{allowed\_signers} \\
Quando escolher & padrão consolidado, usado há muito tempo & mais simples para quem já usa SSH \\
\bottomrule
\end{tabularx}
\end{center}

\begin{dica}
Escolha \textbf{um} dos dois e siga só a seção correspondente. Se não souber qual, comece pelo SSH: são menos passos.
\end{dica}

% -----------------------------------------------------------------------------
\section{Assinando com GPG}

\subsection{Instalar o GnuPG}

\begin{itemize}
  \item \textbf{Linux:} costuma vir instalado. Confira com \comando{gpg --version}. No Debian e no Ubuntu, se faltar: \comando{sudo apt install gnupg}.
  \item \textbf{Windows:} o Git for Windows já traz um \comando{gpg} dentro do \textbf{Git Bash}. Use o Git Bash para todos os comandos desta seção.
\end{itemize}

\begin{aviso}[Windows: um GPG só]
Se você também instalar o Gpg4win, passa a ter \textbf{dois} GPGs no computador, cada um com o próprio chaveiro. Uma chave criada num não aparece no outro, e o Git reclama de \enquote{No secret key}. Crie a chave e assine sempre pelo mesmo. Para saber qual o Git está usando, rode \comando{which gpg} no Git Bash.
\end{aviso}

\subsection{Passo 1 — Ver se você já tem uma chave}

\begin{bashcode}
gpg --list-secret-keys --keyid-format=long
\end{bashcode}

Se aparecer uma chave com o seu e-mail, pule para o passo 3.

\subsection{Passo 2 — Gerar a chave}

\begin{bashcode}
gpg --full-generate-key
\end{bashcode}

O comando faz algumas perguntas:

\begin{center}
\renewcommand{\arraystretch}{1.35}
\small
\begin{tabularx}{\linewidth}{@{}>{\bfseries}l >{\raggedright\arraybackslash}X@{}}
\toprule
\textbf{Pergunta} & \textbf{O que responder} \\
\midrule
Tipo de chave & aperte \tecla{Enter} para aceitar o padrão. Nas versões novas é ECC (ed25519); nas antigas, \enquote{RSA e RSA}. O GitHub aceita os dois \\
Tamanho & só aparece para RSA: digite \texttt{4096} \\
Validade & \texttt{0} = não expira, que é o padrão recomendado pelo GitHub. Uma data (como \texttt{1y}, um ano) é mais segura, mas exige renovar a chave \\
Nome e e-mail & o seu nome e \textbf{o mesmo e-mail verificado} no GitHub \\
Senha & uma senha forte. Ela protege a chave se alguém copiar o seu computador \\
\bottomrule
\end{tabularx}
\end{center}

\subsection{Passo 3 — Descobrir o ID da chave}

\begin{bashcode}
gpg --list-secret-keys --keyid-format=long
\end{bashcode}
\begin{saida}
sec   rsa4096/30F2B65B9246B6CA 2017-08-18 [SC]
      D5E4F29F3275DC0CDA8FFC8730F2B65B9246B6CA
uid                   [ultimate] Maria Exemplo <maria@exemplo.com>
ssb   rsa4096/B7ABC0813E4028C0 2017-08-18 [E]
\end{saida}

O ID é o que vem \textbf{depois da barra} na linha \arq{sec}: aqui, \arq{30F2B65B9246B6CA}. Numa chave ECC, a linha começa com \arq{ed25519/}. O seu ID será outro.

\subsection{Passo 4 — Exportar a chave pública}

\begin{bashcode}
gpg --armor --export 30F2B65B9246B6CA
\end{bashcode}

Copie \textbf{tudo}, da linha \arq{-----BEGIN PGP PUBLIC KEY BLOCK-----} até a linha \arq{-----END PGP PUBLIC KEY BLOCK-----}, incluindo as duas.

\begin{aviso}[Pública, nunca privada]
O que vai para o site é a chave \textbf{pública} (\comando{--export}). Nunca use \comando{--export-secret-keys} para isso. A chave privada não sai do seu computador.
\end{aviso}

\subsection{Passo 5 — Cadastrar no GitHub}

\begin{enumerate}
  \item Clique na sua foto (canto superior direito) \(\rightarrow\)\ \textbf{Settings}.
  \item Na seção \textbf{Access}, clique em \textbf{SSH and GPG keys}.
  \item Clique em \textbf{New GPG key}.
  \item Em \textbf{Title}, dê um nome (por exemplo, \enquote{Notebook pessoal}). Em \textbf{Key}, cole o bloco copiado.
  \item Clique em \textbf{Add GPG key}. O GitHub pode pedir a sua senha de novo.
\end{enumerate}

\subsection{Passo 6 — Dizer ao Git qual chave usar}

\begin{bashcode}
git config --global user.signingkey 30F2B65B9246B6CA
git config --global commit.gpgsign true     # assina todo commit
git config --global tag.gpgSign true        # assina toda tag
\end{bashcode}

\begin{dica}[Já usou assinatura SSH antes?]
Se algum dia configurou \comando{gpg.format ssh}, desfaça com \comando{git config --global --unset gpg.format}. Sem isso, o Git continua tentando assinar com SSH.
\end{dica}

\subsection{Passo 7 — Ensinar o GPG a pedir a senha no terminal}

No Linux e no Git Bash, o GPG precisa saber em qual terminal pedir a senha da chave. Rode uma vez:

\begin{bashcode}
[ -f ~/.bashrc ] && echo -e '\nexport GPG_TTY=$(tty)' >> ~/.bashrc
source ~/.bashrc
\end{bashcode}

Quem usa \textbf{zsh} coloca a linha \arq{export GPG\_TTY=\$(tty)} no \arq{\textasciitilde/.zshrc}.

\subsection{Passo 8 — Testar}

\begin{bashcode}
git commit --allow-empty -m "Testando a assinatura"
git log --show-signature -1
\end{bashcode}
\begin{saida}
commit e55fbf5be5b86eb89737adea712d8c605b5f9529
gpg: Assinatura feita seg 28 set 2026 14:51:16 -03
gpg:                usando RSA chave 0D9FFCDB3C4EBDC3E7225C8CB4BEA86F1F6C7956
gpg: Assinatura correta de "Maria Exemplo <maria@exemplo.com>" [final]
Author: Maria Exemplo <maria@exemplo.com>
\end{saida}

A linha \enquote{Assinatura correta} (em inglês, \enquote{Good signature}) confirma que funcionou. Se você não ligou o \comando{commit.gpgsign}, assine um commit específico com \comando{git commit -S -m "..."}.

% -----------------------------------------------------------------------------
\section{Assinando com SSH}

\subsection{Requisitos}

\begin{itemize}
  \item \textbf{Git 2.34 ou mais novo} (\comando{git --version});
  \item \textbf{OpenSSH 8.1 ou mais novo} (\comando{ssh -V}). A versão 8.7 tem a assinatura quebrada: nela, atualize para a 8.8 ou mais nova;
\end{itemize}

\subsection{Passo 1 — Ter uma chave SSH}

Se você fez o login por SSH no capítulo~\ref{cap:git}, já tem uma: normalmente \arq{\textasciitilde/.ssh/id\_ed25519.pub}. Se não tiver, crie:

\begin{bashcode}
ssh-keygen -t ed25519 -C "maria@exemplo.com"
\end{bashcode}

Aperte \tecla{Enter} para aceitar o local padrão e escolha uma senha. O arquivo terminado em \arq{.pub} é a chave \textbf{pública}. O outro, sem extensão, é a \textbf{privada}, e nunca sai do seu computador.

\subsection{Passo 2 — Dizer ao Git para assinar com SSH}

\begin{bashcode}
git config --global gpg.format ssh
git config --global user.signingkey ~/.ssh/id_ed25519.pub
git config --global commit.gpgsign true
git config --global tag.gpgSign true
\end{bashcode}

\subsection{Passo 3 — Cadastrar como chave de assinatura no GitHub}

\begin{enumerate}
  \item Foto \(\rightarrow\)\ \textbf{Settings} \(\rightarrow\)\ \textbf{SSH and GPG keys} \(\rightarrow\)\ \textbf{New SSH key}.
  \item Em \textbf{Title}, dê um nome. Em \textbf{Key type}, escolha \textbf{Signing Key}.
  \item Em \textbf{Key}, cole o conteúdo do arquivo \arq{.pub} (\comando{cat \textasciitilde/.ssh/id\_ed25519.pub}) e clique em \textbf{Add SSH key}.
\end{enumerate}

\begin{aviso}[No GitHub, a mesma chave entra duas vezes]
Uma chave cadastrada como \emph{Authentication Key} serve para \comando{git push}, mas \textbf{não} para verificar assinaturas. Se quiser usar a mesma chave para as duas coisas, cadastre-a duas vezes: uma com cada tipo.
\end{aviso}

Pelo terminal, com o GitHub CLI, dá no mesmo. A opção \comando{--type} existe nas versões do \comando{gh} lançadas a partir de 2023. Se a sua recusar a opção, use a página web.

\begin{bashcode}
gh ssh-key add ~/.ssh/id_ed25519.pub --type signing --title "Notebook pessoal"
\end{bashcode}

\subsection{Passo 4 — Permitir a conferência no seu computador}

O GitHub confere a assinatura sozinho. Mas, para o \emph{seu} Git conferir, ele precisa de uma lista de chaves confiáveis. Sem ela, aparece:

\begin{saida}
error: gpg.ssh.allowedSignersFile needs to be configured and exist for ssh signature verification
\end{saida}

Crie a lista uma vez:

\begin{bashcode}
touch ~/.ssh/allowed_signers
git config --global gpg.ssh.allowedSignersFile ~/.ssh/allowed_signers
echo "$(git config --get user.email) namespaces=\"git\" $(cat ~/.ssh/id_ed25519.pub)" >> ~/.ssh/allowed_signers
\end{bashcode}

Cada linha do arquivo associa um e-mail a uma chave. O \arq{namespaces="git"} limita a chave a assinaturas do Git.

\subsection{Passo 5 — Testar}

\begin{bashcode}
git commit --allow-empty -m "Testando a assinatura SSH"
git log --show-signature -1
\end{bashcode}
\begin{saida}
commit 86613eb51dee51c80a61f2cf8f1532df427e60c5
Good "git" signature for maria@exemplo.com with ED25519 key SHA256:Z/dW7mOFmUlo/Po1CSObYXoO3bAPY2bTohQtnJY0PGc
Author: Maria Exemplo <maria@exemplo.com>
\end{saida}

% -----------------------------------------------------------------------------
\section{Conferindo no GitHub}

Envie um commit assinado com \comando{git push} e abra a lista de commits:

\begin{itemize}
  \item Na aba \textbf{Commits} do repositório, o selo \textbf{Verified} aparece ao lado do commit. Clique nele para ver qual chave assinou.
\end{itemize}

No terminal, uma letra resume o estado de cada commit:

\begin{bashcode}
git log --format='%h %G? %s'
\end{bashcode}

\begin{center}
\renewcommand{\arraystretch}{1.3}
\begin{tabularx}{0.9\linewidth}{@{}>{\ttfamily\bfseries\color{roxoescuro}}l X@{}}
\toprule
G & assinatura válida \\
N & sem assinatura \\
B & assinatura inválida \\
E & não foi possível conferir (por exemplo, falta a chave pública ou o \arq{allowed\_signers}) \\
\bottomrule
\end{tabularx}
\end{center}

\begin{dica}[Vigilant mode no GitHub]
Em \textbf{Settings \(\rightarrow\)\ SSH and GPG keys}, a opção \textbf{Vigilant mode} (\enquote{Flag unsigned commits as unverified}) faz o GitHub marcar como \textbf{Unverified} todo commit seu que não esteja assinado, inclusive os que alguém fizer se passando por você. Ligue só depois que todos os seus computadores estiverem assinando.
\end{dica}

% -----------------------------------------------------------------------------
\section{Cuidados com a chave}

\begin{itemize}
  \item \textbf{Guarde a chave privada e a senha.} Se perder a chave, é só criar outra e cadastrar, mas perde-se a continuidade.
  \item \textbf{Chave vazou? Remova-a da conta e cadastre outra.} No GitHub, a verificação é registrada no momento do push: os commits antigos \textbf{continuam Verified} mesmo que a chave seja removida, expire ou seja revogada.
  \item \textbf{Commits feitos pela interface web} já saem assinados pela própria plataforma.
  \item No GitHub, um Pull Request integrado com \textbf{Rebase and merge} reescreve os commits, e eles perdem a sua assinatura.
  \item \textbf{Não é preciso publicar a chave em servidor de chaves} (\emph{keyserver}). Alguns tutoriais pedem isso, mas o GitHub só usa a chave cadastrada na conta.
\end{itemize}

\section{Problemas comuns}

\begin{center}
\renewcommand{\arraystretch}{1.4}
\small
\begin{tabularx}{\linewidth}{@{}>{\raggedright\arraybackslash\ttfamily\footnotesize}p{5.6cm} >{\raggedright\arraybackslash}X@{}}
\toprule
\textrm{\textbf{Mensagem ou sintoma}} & \textbf{O que fazer} \\
\midrule
error: gpg failed to sign the data & quase sempre é o \arq{GPG\_TTY} (passo 7 do GPG). Confira também se o \arq{user.signingkey} bate com \comando{gpg --list-secret-keys} \\
gpg: signing failed: Inappropriate ioctl for device & falta o \arq{export GPG\_TTY=\$(tty)} \\
No secret key / secret key not available & o ID configurado não existe neste chaveiro. No Windows, confira se o Git usa o mesmo GPG em que a chave foi criada \\
gpg.ssh.allowedSignersFile needs to be configured and exist for ssh signature verification & crie o \arq{allowed\_signers} (passo 4 do SSH). A assinatura em si está certa \\
either user.signingkey or gpg.ssh.defaultKeyCommand needs to be configured & falta o \comando{git config --global user.signingkey} \\
\textrm{Selo \textbf{Unverified} no site} & o e-mail do commit não está verificado na conta, a chave não foi cadastrada, ou a chave SSH foi cadastrada só como \emph{Authentication Key} \\
\bottomrule
\end{tabularx}
\end{center}

\begin{dica}[Fontes oficiais]
Os passos deste capítulo seguem a documentação do GitHub \cite{github_assinatura}. Os nomes dos botões podem mudar com novas versões do site.
\end{dica}
          % Assinatura de commits (GPG e SSH)
% ======================================================================
\chapter{Exercícios práticos}
\label{cap:exercicios}
% ======================================================================

Faça os exercícios na ordem. Todos os comandos são digitados no \textbf{Git
Bash} (Windows) ou no \textbf{Terminal} (Linux e macOS).

Os materiais dos exercícios (\arq{jogo-da-memoria.zip}, \arq{main.py} e
\arq{modelo-actions.zip}) podem ser baixados pelos botões do topo da página
\url{https://andradasdev.github.io/github/} ou na pasta \arq{materials/} do
repositório \texttt{ronidomingues/github-training}. Os comandos abaixo supõem que eles foram
salvos na pasta \arq{Downloads}.

\begin{dica}[Pasta de treino]
Crie uma pasta só para os exercícios, por exemplo \arq{\textasciitilde/treino},
e entre nela com \comando{cd \textasciitilde/treino} antes de começar.
\end{dica}

% ----------------------------------------------------------------------
\section{Exercício 1 — Conferir as ferramentas}

\begin{bashcode}
git --version
gh --version
gh auth status
\end{bashcode}

\textbf{Como saber se deu certo:} as duas primeiras linhas mostram números de
versão, e a terceira diz \texttt{Logged in to github.com account
<seu-usuario>}. Se não estiver logado, volte ao capítulo~\ref{cap:git}.

% ----------------------------------------------------------------------
\section{Exercício 2 — O primeiro repositório}

\begin{enumerate}
  \item Crie o repositório no GitHub e já traga a cópia para o computador:
\begin{bashcode}
gh repo create meu-primeiro-repo --public --description "Meu primeiro repositório" --clone
cd meu-primeiro-repo
\end{bashcode}
  O Git avisa \texttt{You appear to have cloned an empty repository}. É
  normal: o repositório ainda não tem nenhum arquivo.

  \item Crie o \arq{README.md}, a \enquote{capa} do projeto no GitHub, escrita
        em Markdown:
\begin{bashcode}
cat > README.md << 'EOF'
# Meu Primeiro Projeto

Este é meu primeiro repositório no GitHub!

## O que este projeto faz?

- Aprender Git e GitHub
- Praticar comandos
EOF
\end{bashcode}

  \item Registre e envie:
\begin{bashcode}
git add README.md
git commit -m "docs: cria o README"
git push -u origin main
\end{bashcode}
\end{enumerate}

\begin{conceito}[O que o -u faz]
O \texttt{-u} liga a sua branch \texttt{main} à branch \texttt{main} do
GitHub (\texttt{origin}). Isso só precisa ser feito no primeiro push: depois,
basta \comando{git push} e \comando{git pull}.
\end{conceito}

\textbf{Como saber se deu certo:} a página
\texttt{github.com/<seu-usuario>/meu-primeiro-repo} mostra o README
formatado.

% ----------------------------------------------------------------------
\section{Exercício 3 — Uma licença criada pelo site}

Este exercício mostra o caminho inverso: uma mudança feita no GitHub que
precisa chegar ao seu computador.

\begin{enumerate}
  \item Na página do repositório, clique em \menu{Add file} \(\rightarrow\)
        \menu{Create new file}.
  \item Digite \texttt{LICENSE} como nome. Aparece o botão
        \menu{Choose a license template}: escolha \menu{MIT License},
        confira seu nome e clique em \menu{Review and submit}.
  \item Clique em \menu{Commit changes}.
  \item No terminal, traga a novidade:
\begin{bashcode}
git pull
ls
\end{bashcode}
\end{enumerate}

\textbf{Como saber se deu certo:} o \comando{ls} mostra o arquivo
\arq{LICENSE}, e o \comando{git log --oneline} tem dois commits.

% ----------------------------------------------------------------------
\section{Exercício 4 — Resolver um conflito}

\begin{enumerate}
  \item \textbf{No site:} abra o \arq{README.md} do
        \texttt{meu-primeiro-repo}, clique no lápis (\menu{Edit}), troque a
        primeira linha por \texttt{\# Meu Primeiro Projeto (editado no site)}
        e clique em \menu{Commit changes}.
  \item \textbf{No computador, sem dar pull:} abra o \arq{README.md} num
        editor, troque a \emph{mesma} primeira linha por
        \texttt{\# Meu Primeiro Projeto (editado no computador)}, salve e
        registre:
\begin{bashcode}
git commit -am "docs: edita o título no computador"
git push
\end{bashcode}
  O push é recusado, porque o GitHub tem um commit que você não tem:
\begin{saida}
 ! [rejected]        main -> main (fetch first)
error: failed to push some refs to 'https://github.com/...'
\end{saida}

  \item Traga a alteração do site:
\begin{bashcode}
git pull
\end{bashcode}
\begin{saida}
CONFLICT (content): Merge conflict in README.md
Automatic merge failed; fix conflicts and then commit the result.
\end{saida}

  \item Abra o \arq{README.md}. O Git colocou as duas versões dentro dele:
\begin{saida}
<<<<<<< HEAD
# Meu Primeiro Projeto (editado no computador)
=======
# Meu Primeiro Projeto (editado no site)
>>>>>>> 15261e4...
\end{saida}
  Entre \texttt{<<<<<<<} e \texttt{=======} está a sua versão; entre
  \texttt{=======} e \texttt{>>>>>>>}, a que veio do GitHub. Deixe só a
  linha que deve ficar (ou escreva uma terceira), apague os três marcadores e
  salve.

  \item Registre a resolução e envie:
\begin{bashcode}
git add README.md
git commit -m "docs: resolve o conflito do título"
git push
\end{bashcode}
\end{enumerate}

\textbf{Como saber se deu certo:} \comando{git log --oneline --graph} mostra
as duas linhas de trabalho se juntando no commit de resolução.

\begin{aviso}
Se o \comando{git pull} terminar em \texttt{Need to specify how to reconcile
divergent branches}, rode \comando{git config --global pull.rebase false} e
repita o pull.
\end{aviso}

% ----------------------------------------------------------------------
\section{Exercício 5 — Um site no GitHub Pages}

\begin{enumerate}
  \item Crie o repositório do site:
\begin{bashcode}
cd ~/treino
gh repo create meu-site --public --description "Meu primeiro site" --clone
cd meu-site
\end{bashcode}

  \item Descompacte o \arq{jogo-da-memoria.zip} dentro de \arq{Downloads}
        (no Windows: botão direito \(\rightarrow\) \menu{Extrair tudo}). Ele
        cria a pasta \arq{jogo-da-memoria/}. Copie o \textbf{conteúdo} dela,
        e não a pasta, para dentro do \arq{meu-site}:
\begin{bashcode}
cp -r ~/Downloads/jogo-da-memoria/. .
ls
\end{bashcode}
  O \comando{ls} precisa mostrar o \arq{index.html} aqui, na raiz.

  \item Registre e envie:
\begin{bashcode}
git add .
git commit -m "feat: adiciona o jogo da memória"
git push -u origin main
\end{bashcode}

  \item No GitHub: \menu{Settings} \(\rightarrow\) \menu{Pages}. Em
        \menu{Source}, escolha \menu{Deploy from a branch}; em
        \menu{Branch}, \texttt{main} e \texttt{/(root)}; clique em
        \menu{Save}.
\end{enumerate}

\textbf{Como saber se deu certo:} depois de um ou dois minutos, o jogo abre
em \texttt{https://<seu-usuario>.github.io/meu-site/}.

% ----------------------------------------------------------------------
\section{Exercício 6 — Automação com GitHub Actions}

\begin{enumerate}
  \item Crie o repositório e copie o \arq{main.py} para dentro dele:
\begin{bashcode}
cd ~/treino
gh repo create meu-script-python --public --description "Script Python com GitHub Actions" --clone
cd meu-script-python
cp ~/Downloads/main.py .
\end{bashcode}

  \item Crie a pasta dos workflows e o arquivo:
\begin{bashcode}
mkdir -p .github/workflows
\end{bashcode}
  Abra um arquivo novo \arq{.github/workflows/python.yml} no editor e cole o
  Exemplo 2 do capítulo~\ref{cap:pages}.

  \item Envie:
\begin{bashcode}
git add .
git commit -m "ci: roda o script Python a cada push"
git push -u origin main
\end{bashcode}

  \item Acompanhe na aba \menu{Actions} do repositório, ou com
        \comando{gh run watch}. Quando terminar, traga o arquivo gerado:
\begin{bashcode}
git pull
cat resultado.txt
\end{bashcode}
\end{enumerate}

\textbf{Como saber se deu certo:} o \arq{resultado.txt} mostra o horário da
execução, e o último commit do repositório é do \texttt{github-actions[bot]}.

\begin{dica}[Outro modelo]
O \arq{modelo-actions.zip}, em \arq{materials/}, traz o mesmo fluxo para um
programa em Fortran: compila, executa e grava o resultado.
\end{dica}

% ----------------------------------------------------------------------
\section{Exercício 7 — Fork, Pull Request e lista de presença}
\label{sec:ex_presenca}

A lista de presença deste guia (Anexo~\ref{an:presenca}) é feita por Pull
Request. É o fluxo real de contribuir com o projeto de outra pessoa.

\begin{enumerate}
  \item Crie o seu fork e clone:
\begin{bashcode}
cd ~/treino
gh repo fork ronidomingues/github-training --clone
cd github-training
\end{bashcode}

  \item Crie uma branch para a sua contribuição:
\begin{bashcode}
git switch -c presenca/seu-nome
\end{bashcode}

  \item Crie um arquivo com o seu nome na pasta \arq{presences/}. Use só
        letras, números, espaços, pontos e hífens no nome. Se quiser, escreva
        uma avaliação da capacitação dentro dele.
\begin{bashcode}
echo "Gostei muito da parte de conflitos!" > "presences/Seu Nome.txt"
\end{bashcode}

  \item Registre e envie para o \emph{seu} fork:
\begin{bashcode}
git add presences
git commit -m "chore(presences): registra a presença de Seu Nome"
git push -u origin presenca/seu-nome
\end{bashcode}

  \item Abra o Pull Request. No GitHub, na página do seu fork, clique em
        \menu{Compare \& pull request}, confira que a base é
        \texttt{ronidomingues/github-training} e clique em
        \menu{Create pull request}. Pelo terminal:
\begin{bashcode}
gh pr create --repo ronidomingues/github-training --fill
\end{bashcode}
\end{enumerate}

\textbf{Como saber se deu certo:} o PR aparece na aba \menu{Pull requests} do
repositório original. Depois que ele for aceito, o seu nome entra na lista
de presença da próxima versão do guia.

% ----------------------------------------------------------------------
\section{Exercício 8 — Mensagens de commit}

Reescreva cada mensagem no formato Conventional Commits
(capítulo~\ref{cap:padroes}):

\begin{enumerate}
  \item \enquote{arrumei o erro que dava quando a senha tava vazia}
  \item \enquote{coloquei explicação de como instalar no readme}
  \item \enquote{agora o relatório sai em pdf também}
  \item \enquote{atualizei o workflow pra rodar os testes}
\end{enumerate}

\begin{dica}[Respostas possíveis]
\texttt{fix(login): aceita senha vazia sem erro} ·
\texttt{docs(readme): explica a instalação} ·
\texttt{feat(relatorio): exporta em PDF} ·
\texttt{ci: roda os testes no workflow}
\end{dica}

% ----------------------------------------------------------------------
\section{Exercício 9 — Um commit assinado}

Siga o capítulo~\ref{cap:assinatura} com GPG ou SSH e faça um commit no
\texttt{meu-primeiro-repo}:

\begin{bashcode}
git commit --allow-empty -m "chore: testa a assinatura"
git push
git log --format='%h %G? %s' -1
\end{bashcode}

\textbf{Como saber se deu certo:} o terminal mostra a letra \texttt{G}, e a
aba \menu{Commits} do GitHub mostra o selo \textbf{Verified}.

% ----------------------------------------------------------------------
\section{Checklist}

\begin{itemize}
  \item[$\square$] Git e GitHub CLI instalados e logados
  \item[$\square$] Repositório criado, com README e LICENSE
  \item[$\square$] Um conflito resolvido
  \item[$\square$] Site publicado no GitHub Pages
  \item[$\square$] Workflow rodando no GitHub Actions
  \item[$\square$] Pull Request de presença aberto
  \item[$\square$] Commit assinado com selo \textbf{Verified}
\end{itemize}
         % Exercícios práticos
% ======================================================================
\chapter{Conclusão e próximos passos}
\label{cap:conclusao}
% ======================================================================

Em duas horas, o caminho foi do zero ao fluxo completo de um projeto no
GitHub: conta protegida e benefícios de estudante, Git e GitHub CLI
instalados, commits bem escritos, branches, conflitos resolvidos,
contribuição por fork e Pull Request, um site no ar pelo GitHub Pages, uma
automação no GitHub Actions, arquivos grandes com o Git LFS e commits
assinados.

Mais importante que decorar comandos é entender o modelo por trás deles: o
histórico é uma sequência de commits, uma branch é um ponteiro, e um Pull
Request é uma conversa antes de integrar. Com isso, os comandos novos passam a
fazer sentido sozinhos.

\section{Para continuar estudando}

\begin{itemize}
  \item \textbf{Pro Git}, de Scott Chacon e Ben Straub: o livro de referência
        do Git, gratuito e com tradução para o português,
        \url{https://git-scm.com/book/pt-br/v2}.
  \item \textbf{GitHub Skills}: cursos interativos dentro do próprio GitHub,
        \url{https://skills.github.com}.
  \item \textbf{Learn Git Branching}: exercícios visuais de branch e merge, em
        português, \url{https://learngitbranching.js.org/?locale=pt_BR}.
  \item \textbf{Documentação do GitHub}: \url{https://docs.github.com/pt}.
\end{itemize}

\begin{dica}[O próximo passo]
Escolha um projeto seu, mesmo pequeno, e coloque no GitHub com README,
licença e um workflow. Um repositório bem cuidado vale mais no portfólio do
que muitos repositórios vazios.
\end{dica}
     % Conclusão

% --- Referências -------------------------------------------------------
\nocite{conventional_commits,adorno_commits,iuricode_commits,git_docs}   % fontes consultadas, mesmo sem citação no texto
\bibliographystyle{abntex2-alf}
\bibliography{pages/references}

% --- Elementos pós-textuais -------------------------------------------
% ======================================================================
% Apêndices
% ======================================================================
\begin{apendicesenv}
\partapendices

\chapter{Comandos do Git}
\label{ap:git}

Referência dos principais comandos do Git, organizados por categoria. A
documentação completa está em \url{https://git-scm.com/docs}.

\begin{longtable}{|p{3cm}|p{4cm}|p{7cm}|}
    \caption{Comandos do Git} \label{tab:gitcommands} \\
    \hline
    \textbf{Comando} & \textbf{Categoria} & \textbf{Explicação Detalhada} \\
    \hline
    \endfirsthead
    \multicolumn{3}{c}%
    {\textbf{Continuação da Tabela: Comandos do Git}} \\
    \hline
    \textbf{Comando} & \textbf{Categoria} & \textbf{Explicação Detalhada} \\
    \hline
    \endhead
    \multicolumn{3}{|r|}{Continua na próxima página} \\
    \hline
    \endfoot
    \endlastfoot

    git init & Configuração/Setup & Transforma o diretório atual em um repositório Git, criando o diretório \texttt{.git}. Pode ser executado com segurança em um diretório existente sem sobrescrever configurações. \\
    \hline
    git config & Configuração/Setup & Lê ou define variáveis de configuração em nível de sistema, global ou local. Essencial para definir a identidade (user.name, user.email) do autor do commit. \\
    \hline
    git clone [url] & Configuração/Setup & Cria uma cópia local de um repositório remoto. Configura automaticamente a referência 'origin' e faz o checkout da branch principal. \\
    \hline
    git add [file] & Snapshotting Básico & Move alterações de um arquivo do Working Tree para o Index (Staging Area), preparando-o para o próximo commit. \\
    \hline
    git status & Snapshotting Básico & Exibe o estado da Working Tree e do Index, listando arquivos modificados, staged ou não rastreados. \\
    \hline
    git diff & Snapshotting Básico & Mostra as diferenças entre o Working Tree e o Index (alterações não staged). \\
    \hline
    git diff --staged & Snapshotting Básico & Mostra as diferenças entre o Index (Staging Area) e o último commit (HEAD). \\
    \hline
    git commit -m "[msg]" & Snapshotting Básico & Salva o conteúdo atualmente no Index como um novo snapshot permanente (commit) na história. \\
    \hline
    git commit --amend & Manipulação Histórico & Altera o commit anterior, seja modificando sua mensagem ou adicionando/removendo arquivos. Isso reescreve o histórico, gerando um novo SHA. \\
    \hline
    git rm [file] & Gerenciamento Arquivo & Remove um arquivo do Working Tree e do Index. O uso de \texttt{--cached} remove apenas do Index, mantendo o arquivo local. \\
    \hline
    git mv [old][new] & Gerenciamento Arquivo & Move ou renomeia um arquivo de forma rastreada pelo Git. \\
    \hline
    git clean & Gerenciamento Arquivo & Remove arquivos não rastreados (untracked files) do Working Tree. \\
    \hline
    git reset --soft [hash] & Manipulação Histórico & Move o ponteiro HEAD para o commit, mas mantém o Index e o Working Tree intactos (alterações permanecem staged). \\
    \hline
    git reset --mixed [hash] & Manipulação Histórico & (Padrão) Move o HEAD para o commit e reseta o Index (desencena arquivos), preservando o Working Tree. \\
    \hline
    git reset --hard [hash] & Manipulação Histórico & Move o HEAD e reseta o Index e o Working Tree, descartando todas as mudanças locais desde o hash. Altamente destrutivo. \\
    \hline
    git branch & Branching/Navegação & Gerenciamento de branches: lista, cria ou deleta branches locais. \\
    \hline
    git checkout & Branching/Navegação & Comando legado multi-uso. Alterna entre branches ou restaura arquivos antigos/commits, podendo resultar em 'detached HEAD'. \\
    \hline
    git switch & Branching/Navegação & Comando moderno focado em alternar branches. Atualiza a Working Tree e o Index. Utilizado para criar novas branches de forma segura. \\
    \hline
    git merge [branch] & Integração/Merge & Integra alterações de uma branch na atual, criando um 'merge commit' se houver divergência. Operação não-destrutiva. \\
    \hline
    git rebase [base] & Integração/Rebase & Move ou reaplica commits para uma nova base, reescrevendo o histórico para mantê-lo linear. Ideal para branches locais e não publicadas. \\
    \hline
    git rebase -i [base] & Integração/Rebase & Modo interativo do rebase, permitindo squash (combinação), edição ou reordenação de commits. \\
    \hline
    git cherry-pick [hash] & Integração/Portabilidade & Aplica as alterações introduzidas por um único commit específico na branch atual, criando um novo commit equivalente. \\
    \hline
    git revert [hash] & Manipulação Histórico & Cria um novo commit que desfaz as alterações introduzidas por um commit anterior. Usado para desfazer mudanças em histórico compartilhado de forma segura. \\
    \hline
    git fetch & Sincronização Remota & Baixa dados (objetos e refs) de um repositório remoto para o repositório local, sem alterar o Working Tree ou Index (operação segura). \\
    \hline
    git pull & Sincronização Remota & Equivalente a git fetch seguido por uma integração (default: merge). Pode alterar o estado local e causar conflitos imediatamente (operação menos segura). \\
    \hline
    git push [remote][branch] & Sincronização Remota & Carrega commits locais para um repositório remoto. Exige uma operação fast-forward, a menos que \texttt{--force} seja utilizado. \\
    \hline
    git push --tags & Sincronização Remota & Envia tags locais para o repositório remoto. \\
    \hline
    git remote & Sincronização Remota & Gerencia os repositórios remotos rastreados (e.g., listar, adicionar, remover). \\
    \hline
    git log & Auditoria/Inspeção & Exibe o histórico de commits. \\
    \hline
    git shortlog & Auditoria/Inspeção & Fornece um resumo conciso do git log, agrupando commits por autor. \\
    \hline
    git show & Auditoria/Inspeção & Exibe informações detalhadas sobre um objeto Git (commit, tag, etc.). \\
    \hline
    git reflog & Auditoria/Recuperação & Registra as atualizações locais no HEAD e em outras referências, agindo como uma rede de segurança para recuperar commits perdidos após resets ou rebase. \\
    \hline
    git tag & Marcação/Utilitários & Cria, lista, deleta ou verifica objetos de tag, usados para marcar pontos estáticos (releases) no histórico. \\
    \hline
    git tag -a [name] & Marcação/Utilitários & Cria uma tag anotada (com metadados e mensagem), preferida para releases públicas. \\
    \hline
    git stash & Utilitários de Contexto & Salva temporariamente o Working Directory e o Index (alterações ainda sem commit) para permitir a troca de contexto. \\
    \hline
    git stash pop & Utilitários de Contexto & Aplica o último stash salvo e o remove da lista de stashes. \\
    \hline
    git stash apply & Utilitários de Contexto & Aplica o último stash salvo, mas o mantém na lista. \\
    \hline
    git submodule & Utilitários Avançados & Inicializa, atualiza ou inspeciona submódulos (repositórios aninhados). \\
    \hline
    git worktree & Utilitários Avançados & Gerencia múltiplas Working Trees (checkouts) do mesmo repositório, permitindo trabalhar em várias branches simultaneamente. \\
    \hline
    gitk & Utilitários Avançados & O navegador de repositório Git (ferramenta GUI). \\
    \hline
    scalar & Utilitários Avançados & Ferramenta projetada para gerenciar repositórios Git de grande escala (Large Git Repositories). \\
    \hline
    git sparse-checkout & Utilitários Avançados & Reduz a Working Tree para um subconjunto de arquivos rastreados, otimizando o desempenho em repositórios massivos. \\
    \hline
\end{longtable}
\chapter{Comandos do GitHub CLI}
\label{ap:gh}

Referência dos principais comandos do GitHub CLI (\texttt{gh}), organizados
por grupo. A lista completa, sempre atualizada, está no manual oficial
\cite{gh_manual}.

\begin{longtable}{|p{2cm}|p{2.3cm}|p{4cm}|p{6cm}|}
    \caption{Comandos do GitHub CLI (gh)} \label{tab:githubcli_commands} \\
    \hline
    \textbf{Comando Base} & \textbf{Subcomando} & \textbf{Explicação Funcional Detalhada} & \textbf{Exemplo de Sintaxe Chave} \\
    \hline
    \endfirsthead
    \multicolumn{4}{c}%
    {\textbf{Continuação da Tabela: Comandos do GitHub CLI (gh)}} \\
    \hline
    \textbf{Comando Base} & \textbf{Subcomando} & \textbf{Explicação Funcional Detalhada} & \textbf{Exemplo de Sintaxe Chave} \\
    \hline
    \endhead
    \multicolumn{4}{|r|}{\textit{Continua na próxima página}} \\
    \hline
    \endfoot
    \endlastfoot
    
    gh alias & set & Cria um alias para um comando gh, permitindo atalhos personalizados para comandos frequentes & \texttt{gh alias set prc "pr create"} \\
    \hline
    gh alias & list & Lista todos os aliases configurados no GitHub CLI & \texttt{gh alias list} \\
    \hline
    gh alias & delete & Remove um alias previamente configurado & \texttt{gh alias delete prc} \\
    \hline
    gh auth & login & Autentica o usuário no GitHub via navegador web ou token & \texttt{gh auth login} \\
    \hline
    gh auth & logout & Remove a autenticação do usuário atual & \texttt{gh auth logout} \\
    \hline
    gh auth & status & Exibe o status de autenticação atual e usuário conectado & \texttt{gh auth status} \\
    \hline
    gh auth & refresh & Renova a autenticação para um host específico & \texttt{gh auth refresh -{-}hostname github.com} \\
    \hline
    gh auth & token & Exibe o token de autenticação atual & \texttt{gh auth token} \\
    \hline
    gh browse & - & Abre o repositório atual no navegador web & \texttt{gh browse} \\
    \hline
    gh browse & -{-}branch & Abre uma branch específica no navegador & \texttt{gh browse -{-}branch feature-branch} \\
    \hline
    gh browse & -{-}commit & Abre um commit específico no navegador & \texttt{gh browse -{-}commit abc123} \\
    \hline
    gh browse & <número> & Abre a issue ou o pull request de número informado no navegador & \texttt{gh browse 42} \\
    \hline
    gh browse & -{-}settings & Abre as configurações do repositório no navegador & \texttt{gh browse -{-}settings} \\
    \hline
    gh browse & -{-}wiki & Abre a wiki do repositório no navegador & \texttt{gh browse -{-}wiki} \\
    \hline
    gh codespace & code & Abre um codespace no Visual Studio Code & \texttt{gh codespace code} \\
    \hline
    gh codespace & cp & Copia arquivos entre o sistema local e um codespace & \texttt{gh codespace cp local.txt remote:./} \\
    \hline
    gh codespace & create & Cria um novo codespace & \texttt{gh codespace create} \\
    \hline
    gh codespace & delete & Remove um codespace específico & \texttt{gh codespace delete -c my-codespace} \\
    \hline
    gh codespace & jupyter & Abre um codespace no JupyterLab & \texttt{gh codespace jupyter} \\
    \hline
    gh codespace & list & Lista todos os codespaces disponíveis & \texttt{gh codespace list} \\
    \hline
    gh codespace & logs & Exibe os logs de um codespace específico & \texttt{gh codespace logs -c my-codespace} \\
    \hline
    gh codespace & ports & Lista as portas encaminhadas de um codespace & \texttt{gh codespace ports} \\
    \hline
    gh codespace & ports forward & Encaminha uma porta do codespace para o local & \texttt{gh codespace ports forward 3000:4000} \\
    \hline
    gh codespace & ports visibility & Define a visibilidade de uma porta & \texttt{gh codespace ports visibility 3000:public} \\
    \hline
    gh codespace & ssh & Conecta-se a um codespace via SSH & \texttt{gh codespace ssh} \\
    \hline
    gh codespace & stop & Para um codespace em execução & \texttt{gh codespace stop -c my-codespace} \\
    \hline
    gh gist & create & Cria um novo gist a partir de arquivos ou entrada padrão & \texttt{gh gist create script.py} \\
    \hline
    gh gist & clone & Clona um gist específico para o sistema local & \texttt{gh gist clone abc123} \\
    \hline
    gh gist & delete & Remove um gist específico & \texttt{gh gist delete abc123} \\
    \hline
    gh gist & edit & Edita um gist existente & \texttt{gh gist edit abc123} \\
    \hline
    gh gist & list & Lista todos os gists do usuário & \texttt{gh gist list} \\
    \hline
    gh gist & view & Visualiza um gist específico no terminal & \texttt{gh gist view abc123} \\
    \hline
    gh issue & create & Cria uma nova issue no repositório & \texttt{gh issue create -{-}title "Bug" -{-}body "Descrição"} \\
    \hline
    gh issue & list & Lista issues do repositório com filtros opcionais & \texttt{gh issue list -{-}state open} \\
    \hline
    gh issue & status & Mostra o status das issues relevantes para o usuário & \texttt{gh issue status} \\
    \hline
    gh issue & close & Fecha uma issue específica & \texttt{gh issue close 42} \\
    \hline
    gh issue & comment & Adiciona um comentário a uma issue & \texttt{gh issue comment 42 -{-}body "Comentário"} \\
    \hline
    gh issue & delete & Remove uma issue específica & \texttt{gh issue delete 42} \\
    \hline
    gh issue & edit & Edita uma issue existente & \texttt{gh issue edit 42 -{-}title "Novo título"} \\
    \hline
    gh issue & lock & Trava os comentários de uma issue & \texttt{gh issue lock 42} \\
    \hline
    gh issue & reopen & Reabre uma issue fechada & \texttt{gh issue reopen 42} \\
    \hline
    gh issue & transfer & Transfere uma issue para outro repositório & \texttt{gh issue transfer 42 owner/repo} \\
    \hline
    gh issue & view & Exibe detalhes de uma issue específica & \texttt{gh issue view 42} \\
    \hline
    gh project & copy & Copia um projeto para outro dono (usuário ou organização) & \texttt{gh project copy 1 -{-}source-owner @me -{-}target-owner novaorg -{-}title "Cópia" -{-}drafts} \\
    \hline
    gh project & create & Cria um novo projeto & \texttt{gh project create -{-}owner @me -{-}title "Meu Projeto"} \\
    \hline
    gh project & delete & Remove um projeto específico & \texttt{gh project delete 1 -{-}owner @me} \\
    \hline
    gh project & edit & Edita as propriedades de um projeto & \texttt{gh project edit 1 -{-}owner @me -{-}title "Novo Título"} \\
    \hline
    gh project & field-create & Cria um campo personalizado no projeto & \texttt{gh project field-create 1 -{-}owner @me -{-}name "Prioridade" -{-}data-type TEXT} \\
    \hline
    gh project & item-add & Adiciona uma issue ou um pull request ao projeto & \texttt{gh project item-add 1 -{-}owner @me -{-}url \url{https://github.com/owner/repo/issues/1}} \\
    \hline
    gh project & list & Lista projetos disponíveis & \texttt{gh project list -{-}owner owner} \\
    \hline
    gh project & view & Visualiza detalhes de um projeto específico & \texttt{gh project view 1 -{-}owner @me} \\
    \hline
    gh pr & checks & Exibe os status checks de um pull request & \texttt{gh pr checks 15} \\
    \hline
    gh pr & close & Fecha um pull request específico & \texttt{gh pr close 15} \\
    \hline
    gh pr & comment & Adiciona um comentário a um pull request & \texttt{gh pr comment 15 -{-}body "Comentário"} \\
    \hline
    gh pr & create & Cria um novo pull request & \texttt{gh pr create -{-}title "Feature" -{-}body "Descrição"} \\
    \hline
    gh pr & diff & Exibe as diferenças introduzidas pelo pull request & \texttt{gh pr diff 15} \\
    \hline
    gh pr & edit & Edita propriedades de um pull request & \texttt{gh pr edit 15 -{-}title "Novo Título"} \\
    \hline
    gh pr & list & Lista pull requests do repositório & \texttt{gh pr list -{-}state open} \\
    \hline
    gh pr & merge & Mescla um pull request & \texttt{gh pr merge 15 -{-}squash} \\
    \hline
    gh pr & ready & Marca um pull request como pronto para revisão & \texttt{gh pr ready 15} \\
    \hline
    gh pr & reopen & Reabre um pull request fechado & \texttt{gh pr reopen 15} \\
    \hline
    gh pr & review & Adiciona uma revisão a um pull request & \texttt{gh pr review 15 -{-}approve} \\
    \hline
    gh pr & status & Mostra o status dos pull requests relevantes & \texttt{gh pr status} \\
    \hline
    gh pr & view & Exibe detalhes de um pull request específico & \texttt{gh pr view 15} \\
    \hline
    gh pr & checkout & Faz checkout da branch de um pull request & \texttt{gh pr checkout 15} \\
    \hline
    gh release & create & Cria um novo release & \texttt{gh release create v1.0.0 -{-}title "Versão 1.0.0"} \\
    \hline
    gh release & delete & Remove um release específico & \texttt{gh release delete v1.0.0} \\
    \hline
    gh release & download & Baixa os assets de um release & \texttt{gh release download v1.0.0} \\
    \hline
    gh release & list & Lista todos os releases do repositório & \texttt{gh release list} \\
    \hline
    gh release & upload & Faz upload de assets para um release & \texttt{gh release upload v1.0.0 arquivo.zip} \\
    \hline
    gh release & view & Exibe detalhes de um release específico & \texttt{gh release view v1.0.0} \\
    \hline
    gh release & edit & Edita propriedades de um release existente & \texttt{gh release edit v1.0.0 -{-}title "Novo Título"} \\
    \hline
    gh repo & archive & Arquiva um repositório & \texttt{gh repo archive owner/repo} \\
    \hline
    gh repo & clone & Clona um repositório para o sistema local & \texttt{gh repo clone owner/repo} \\
    \hline
    gh repo & create & Cria um novo repositório & \texttt{gh repo create meu-repo -{-}public} \\
    \hline
    gh repo & delete & Remove um repositório & \texttt{gh repo delete owner/repo} \\
    \hline
    gh repo & edit & Edita propriedades de um repositório & \texttt{gh repo edit -{-}description "Nova descrição"} \\
    \hline
    gh repo & fork & Cria um fork de um repositório & \texttt{gh repo fork owner/repo} \\
    \hline
    gh repo & list & Lista repositórios do usuário ou organização & \texttt{gh repo list -{-}limit 10} \\
    \hline
    gh repo & rename & Renomeia um repositório & \texttt{gh repo rename novo-nome} \\
    \hline
    gh repo & sync & Sincroniza um fork com seu repositório upstream & \texttt{gh repo sync} \\
    \hline
    gh repo & view & Exibe detalhes de um repositório & \texttt{gh repo view owner/repo} \\
    \hline
    gh repo & deploy-key & Gerencia chaves de deploy do repositório & \texttt{gh repo deploy-key add chave.pub -{-}title "Servidor"} \\
    \hline
    gh run & cancel & Cancela uma execução de workflow & \texttt{gh run cancel 123456789} \\
    \hline
    gh run & delete & Remove execuções de workflow & \texttt{gh run delete 123456789} \\
    \hline
    gh run & download & Baixa artifacts de uma execução & \texttt{gh run download 123456789} \\
    \hline
    gh run & list & Lista execuções de workflows & \texttt{gh run list} \\
    \hline
    gh run & rerun & Reexecuta um workflow falho & \texttt{gh run rerun 123456789} \\
    \hline
    gh run & view & Exibe detalhes de uma execução & \texttt{gh run view 123456789} \\
    \hline
    gh run & watch & Monitora uma execução em tempo real & \texttt{gh run watch 123456789} \\
    \hline
    gh search & code & Busca por código no GitHub & \texttt{gh search code "função javascript"} \\
    \hline
    gh search & commits & Busca por commits & \texttt{gh search commits "fix bug" -{-}author=user} \\
    \hline
    gh search & issues & Busca por issues e pull requests & \texttt{gh search issues "bug label:bug"} \\
    \hline
    gh search & prs & Busca especificamente por pull requests & \texttt{gh search prs "feature state:open"} \\
    \hline
    gh search & repos & Busca por repositórios & \parbox{3cm}{\texttt{gh search repos "topic:machine-learning"}} \\
    \hline
    gh search & users & Busca por usuários & \texttt{gh search users "nome location:Brasil"} \\
    \hline
    gh secret & list & Lista secrets disponíveis & \texttt{gh secret list} \\
    \hline
    gh secret & remove & Remove um secret específico & \texttt{gh secret remove API\textunderscore KEY} \\
    \hline
    gh secret & set & Define ou atualiza um secret & \texttt{gh secret set API\textunderscore KEY -{-}body "valor"} \\
    \hline
    gh ssh-key & add & Adiciona uma chave SSH à conta & \texttt{gh ssh-key add chave.pub -{-}title "Laptop"} \\
    \hline
    gh ssh-key & list & Lista chaves SSH da conta & \texttt{gh ssh-key list} \\
    \hline
    gh ssh-key & delete & Remove uma chave SSH & \texttt{gh ssh-key delete 123} \\
    \hline
    gh workflow & disable & Desabilita um workflow & \texttt{gh workflow disable "CI Tests"} \\
    \hline
    gh workflow & enable & Habilita um workflow & \texttt{gh workflow enable "CI Tests"} \\
    \hline
    gh workflow & list & Lista workflows disponíveis & \texttt{gh workflow list} \\
    \hline
    gh workflow & run & Executa um workflow manualmente & \texttt{gh workflow run "CI Tests"} \\
    \hline
    gh workflow & view & Exibe detalhes de um workflow & \texttt{gh workflow view "CI Tests"} \\
    \hline
\end{longtable}

% ----------------------------------------------------------------------
\chapter{Glossário}
\label{ap:glossario}

\begin{description}
  \item[Branch] Linha de desenvolvimento dentro do histórico; tecnicamente, um
        ponteiro para um commit.
  \item[Clone] Cópia completa de um repositório remoto no seu computador.
  \item[Commit] Registro de um estado do projeto, com autor, data e mensagem.
  \item[Conflito] Situação em que duas alterações na mesma região de um
        arquivo não podem ser combinadas automaticamente.
  \item[Fork] Cópia de um repositório de outra pessoa na sua conta, para você
        propor mudanças por Pull Request.
  \item[GitHub Actions] Plataforma de automação do GitHub, descrita em
        arquivos YAML.
  \item[GitHub Pages] Hospedagem de sites estáticos a partir de um
        repositório.
  \item[.gitignore] Arquivo que lista o que o Git não deve versionar.
  \item[Merge] Integração de duas linhas de histórico.
  \item[Origin] Nome padrão do repositório remoto principal.
  \item[Pull Request] Pedido para integrar uma branch em outra, com revisão.
  \item[README] Arquivo de apresentação do repositório, em Markdown.
  \item[Remoto] Cópia do repositório em outro lugar, como o GitHub.
  \item[Repositório] Pasta cujo histórico é controlado pelo Git.
  \item[Workflow] Processo automatizado do GitHub Actions.
\end{description}

\end{apendicesenv}

% ======================================================================
% Anexos
% ======================================================================
\begin{anexosenv}
\partanexos

\chapter{Lista de presença}
\label{an:presenca}

Participantes que registraram a presença por Pull Request
(Exercício~7, seção~\ref{sec:ex_presenca}). A data é a do momento em que o
Pull Request foi integrado ao repositório.

% O arquivo abaixo é gerado por scripts/generate_presence_list.py a partir da
% pasta presences/. Sem ele (compilação local), o guia avisa em vez de falhar.
\IfFileExists{pages/presence_list.tex}{%
  \begin{longtable}{@{}p{0.6\textwidth} p{0.35\textwidth}@{}}
\toprule
\textbf{Nome} & \textbf{Presença registrada em} \\
\midrule
\endhead
Ronivaldo D. Andrade & 09/10/2025 às 12:04 \\
Persona Presente & 09/10/2025 às 14:35 \\
\bottomrule
\end{longtable}
%
}{%
  \textit{Lista gerada automaticamente na compilação do repositório. Rode
  \texttt{python3 scripts/generate\_presence\_list.py} para gerá-la localmente.}%
}

\end{anexosenv}


\end{document}
